\documentclass[12pt,a4paper]{article}
\usepackage{fontspec}
\setmainfont{DejaVu Serif}
\usepackage[margin=2.2cm]{geometry}
\usepackage{amsmath,amssymb,amsfonts}
\usepackage{graphicx}
\usepackage{booktabs}
\usepackage{tabularx}
\usepackage{array}
\usepackage{float}
\usepackage{xcolor}
\usepackage{caption}
\usepackage{listings}
\definecolor{codebg}{RGB}{241,245,249}
\lstset{frame=single,backgroundcolor=\color{codebg},basicstyle=\ttfamily\small,breaklines=true,breakatwhitespace=true,linewidth=\linewidth}

\usepackage{hyperref}

\hypersetup{
    colorlinks=true,
    breaklinks=true,
    linkcolor=blue,
    filecolor=magenta,      
    urlcolor=cyan,
    pdftitle={Bao cao Phan 6: Thiet ke, Huan luyen va Danh gia Mo hinh Du bao San luong Quang dien (Khong Tran Le Phai)}
}

\begin{document}

\begin{titlepage}
    \centering
    \vspace*{1cm}
    
    {\Huge \bfseries BÁO CÁO PHẦN 6 \par}
    \vspace{0.5cm}
    {\huge \bfseries THIẾT KẾ, HUẤN LUYỆN VÀ ĐÁNH GIÁ MÔ HÌNH DỰ BÁO SẢN LƯỢNG QUANG ĐIỆN \par}
    \vspace{0.3cm}
    {\large \bfseries BỔ SUNG CÁC KỸ THUẬT MÃ HÓA, CHUẨN HÓA \& CODE TỪ NOTEBOOKS \par}
    \vspace{0.2cm}
    {\large \itshape (Solar Forecasting Model Development \& Notebook Encoding Code Snippets) \par}
    
    \vspace{1.5cm}
    \textbf{Tác giả:} Nhóm Nghiên Cứu "The Outliers" --- Đồ Án Tốt Nghiệp \\
    \textbf{Đối tượng nghiên cứu:} Hệ thống dự báo sản lượng điện mặt trời cho 42 trạm quang điện độc lập \\
    \textbf{Tần suất dự báo:} 15 phút ($h=1$: 15m, $h=4$: 1h) \\
    \textbf{Trạng thái:} Đã nghiệm thu \& kiểm chứng độc lập trên tập Test Niêm Phong (Sealed Test Set)
    
    \vfill
    {\large TP. Hồ Chí Minh --- NĂM 2026 \par}
\end{titlepage

\end{titlepage}

\begin{center}
\begin{minipage}{0.92\linewidth}
\small
\textbf{\large TÓM TẮT (ABSTRACT)} \vspace{0.2cm} \\
Dự báo sản lượng điện mặt trời theo từng chu kỳ 15 phút ($h=1$) giữ vai trò quyết định trong việc điều tiết hệ thống điện, duy trì công suất quay dự phòng và giảm thiểu phạt trệch cân bằng trên thị trường năng lượng. Báo cáo này trình bày nghiên cứu toàn diện về việc xây dựng, tối ưu hóa và đánh giá mô hình dự báo chuỗi thời gian cho 42 trang trại quang điện độc lập. Nghiên cứu đóng góp 3 giải pháp kỹ thuật cốt lõi: (1) Thuật toán phối hợp lọc bất thường kép GMM-IF (Gaussian Mixture Model \& Isolation Forest Fusion, chuẩn bài báo \textit{IEEE Access 2025}) kết hợp rào chắn kiểm soát trần công suất vật lý ($P_{\text{capacity}} \times 1.02$); (2) Phương pháp triệt tiêu hoàn toàn hiện tượng lệch pha (Phase Lag Elimination) bằng cách cấm sử dụng đặc trưng sao chép $lag\_1$ ($t-15\text{m}$) và tái tạo bức xạ mặt trời theo mô hình Haurwitz Clear-Sky; và (3) Hệ thống Executive BI Dashboard chuyên nghiệp trên nền tảng Streamlit với 17 biểu đồ giám sát thực thời và giải thích minh bạch bằng trị số SHAP (Nobel Kinh tế 2012). Kết quả thực nghiệm kiểm chứng độc lập trên Tập Test Niêm Phong ($15\%$ dữ liệu) khẳng định mô hình đề xuất đạt chỉ số WAPE ban ngày **17.18\%** (h=1) và **20.75\%** (h=4), MAE **1.2920 kWh** (h=1), $R^2$ **0.9288** (h=1), và Forecast Skill Score **+71.83\%** so với Prophet Baseline (h=1), đồng thời giữ độ trễ pha tuyệt đối **+0.00 phút** trên toàn bộ 42 trạm. So với baseline Persistence (đo trên tập Validation), mô hình vượt trội ở cả 2 horizon, rõ rệt hơn tại horizon xa (h=1: **+15.1\%**, h=4: **+49.3\%**, xem mục Benchmark chi tiết). \vspace{0.3cm} \\
\textbf{Từ khóa (Keywords):} Dự báo quang điện 15 phút, LightGBM, Khử trễ pha (Phase Lag Elimination), Outlier Fusion GMM-IF, SHAP (XAI), Executive BI Dashboard.
\end{minipage}
\end{center}
\vspace{0.5cm}
}

\tableofcontents
\newpage

\section[TỔNG QUAN VÀ BỐI CẢNH NGHIỆP VỤ]{TỔNG QUAN VÀ BỐI CẢNH NGHIỆP VỤ\\(BUSINESS CONTEXT \& GOALS)}

Điện mặt trời là nguồn năng lượng tái tạo có tính \textbf{biến động cao (intermittent)} và phụ thuộc chặt chẽ vào các điều kiện khí tượng như bức xạ mặt trời toàn phần ($GHI$), bức xạ trực tiếp ($DNI$), độ mây phủ ($cloud\_cover$) và nhiệt độ môi trường. Trong thực tế ngành điện, việc dự báo sản lượng quang điện với tần suất 15 phút không chỉ là một bài toán kỹ thuật dữ liệu (Data Science), mà mang \textbf{ý nghĩa nghiệp vụ tài chính và vận hành hạ tầng lưới điện cực kỳ quan trọng}.

\subsection{Ý nghĩa nghiệp vụ đối với Trung tâm Điều độ Hệ thống Điện (A0 / NSMO)}
\begin{enumerate}
    \item \textbf{Lập lịch phát điện tối ưu (Unit Commitment \& Economic Dispatch):} Dự báo chính xác sản lượng điện mặt trời theo từng block 15 phút giúp đơn vị điều độ chủ động điều tiết các nguồn điện truyền thống (thủy điện tích năng, tua-bin khí), giảm chi phí duy trì công suất quay dự phòng (\textit{Spinning Reserve Cost}).
    \item \textbf{Đảm bảo an toàn tần số và điện áp lưới (Grid Frequency \& Voltage Stability):} Việc điện mặt trời sụt giảm chớp nhoáng do hiện tượng mây phủ (\textit{Cloud Transient Event}) nếu không được dự báo trước sẽ gây ra biến động tần số, cơ nguy quá tải cục bộ trên đường dây truyền tải.
\end{enumerate}

\subsection{Ý nghĩa nghiệp vụ đối với Chủ đầu tư trang trại quang điện}
\begin{enumerate}
    \item \textbf{Tối ưu hóa doanh thu \& Giảm chi phí phạt trệch cân bằng (Imbalance Settlement Penalties):} Theo quy định thị trường điện, các nhà máy điện mặt trời phải nộp bản tin dự báo công suất phát trước mỗi chu kỳ giao dịch. Dự báo sai lệch vượt quá ngưỡng cho phép sẽ chịu mức phí phạt trệch cân bằng công suất rất nặng.
    \item \textbf{Giám sát sức khỏe vật lý tài sản (Performance Ratio - PR \& Capacity Factor - CF):} Mô hình dự báo đóng vai trò là "mức cơ sở lý thuyết chuẩn" để đối chiếu với sản lượng thực đo. Nếu sản lượng thực tế thấp hơn hẳn dự báo trong điều kiện bức xạ cao, mô hình giúp phát hiện sớm hiện tượng suy giảm hiệu suất tấm pin (\textit{Panel Degradation}), quá nhiệt Inverter hoặc bám bẩn bề mặt (\textit{Soiling Loss}).
\end{enumerate}


\subsection{Danh Mục Thuật Ngữ \& Từ Viết Tắt Khoa Học (Glossary of Terms \& Acronyms)}
Để bảo đảm tính rõ ràng, minh bạch và dễ tra cứu cho Hội đồng Đánh giá Đồ án, Bảng \ref{tab:glossary} giải thích chi tiết ý nghĩa kỹ thuật, công thức và chuẩn quốc tế của toàn bộ các từ viết tắt được sử dụng trong báo cáo:

\begin{table}[H]
\centering
\caption{Bảng Tra Cứu Chi Tiết Các Thuật Ngữ Kỹ Thuật \& Từ Viết Tắt Khoa Học trong Báo Cáo}
\label{tab:glossary}
\small
\begin{tabularx}{\linewidth}{|l|l|X|}
\hline
\textbf{Ký Hiệu} & \textbf{Tên Gốc / Tiếng Anh} & \textbf{Giải Thích Ý Nghĩa Kỹ Thuật \& Chuẩn Quốc Tế} \\ \hline
\textbf{STC} & Standard Test Conditions & \textbf{Điều kiện Thử nghiệm Tiêu chuẩn} cho tấm pin mặt trời theo chuẩn quốc tế \textbf{IEC 60904-3} (gồm: Bức xạ $1000 \text{ W/m}^2$, Nhiệt độ cell pin $25^\circ\text{C}$, và Khối không khí $\text{AM} = 1.5$). $P_{\text{stc}}$ là công suất định mức ($\text{kWp}$) do nhà sản xuất chứng nhận ở STC. \\ \hline
\textbf{GHI} & Global Horizontal Irradiance & \textbf{Tổng Bức xạ Mặt trời trên Mặt phẳng Ngang} ($\text{W/m}^2$), gồm tổng của Bức xạ Trực tiếp ($DNI$) và Bức xạ Khuếch tán ($DHI$). \\ \hline
\textbf{DNI} & Direct Normal Irradiance & \textbf{Bức xạ Trực tiếp Vuông góc} với tia nắng mặt trời ($\text{W/m}^2$), chiếu thẳng từ đĩa mặt trời khi không bị mây che. \\ \hline
\textbf{DHI} & Diffuse Horizontal Irradiance & \textbf{Bức xạ Khuếch tán Bầu trời} trên mặt phẳng ngang ($\text{W/m}^2$), tạo ra do ánh sáng mặt trời bị tán xạ bởi mây, bụi và khí quyển. \\ \hline
\textbf{POA} & Plane of Array Irradiance & \textbf{Bức xạ trên Mặt phẳng Nghiêng Dàn Pin} ($\text{W/m}^2$), là tổng bức xạ thực tế chiếu lên bề mặt tấm pin lắp nghiêng. \\ \hline
\textbf{PR} & Performance Ratio & \textbf{Tỷ số Hiệu suất Hệ thống} (\%), đo tỷ lệ giữa năng lượng phát ra thực tế so với năng lượng lý thuyết thu nhận ở điều kiện bức xạ $POA$. \\ \hline
\textbf{CF} & Capacity Factor & \textbf{Hệ số Sử dụng Công suất} (\%), tỷ lệ giữa sản lượng phát thực tế trong ngày so với sản lượng tối đa nếu phát liên tục ở $P_{\text{stc}}$ trong 24h. \\ \hline
\textbf{GMM} & Gaussian Mixture Model & \textbf{Mô hình Hỗn hợp Gauss}: Thuật toán học máy thống kê biểu diễn phân bố dữ liệu bằng tổng các hàm mật độ Gauss, dùng để tính mật độ xác suất hậu phương $p_{\text{gmm}} < \theta$ lọc Outlier cục bộ. \\ \hline
\textbf{IF / iForest} & Isolation Forest & \textbf{Rừng Cô lập}: Thuật toán phân lập điểm bất thường dựa trên cấu trúc cây phân tách ngẫu nhiên. Số bước cắt ngắn $h(x)$ thể hiện mức độ dị biệt $S(x, n) > 0.65$ toàn cục (\textit{IEEE Access 2025}). \\ \hline
\textbf{VIF} & Variance Inflation Factor & \textbf{Hệ số Phóng đại Phương sai}: Chỉ số đo lường mức độ đa cộng tuyến giữa các đặc trưng đầu vào. $VIF > 10$ cho thấy đặc trưng bị trùng lặp thông tin nghiêm trọng. \\ \hline
\textbf{SHAP} & SHapley Additive exPlanations & Phương pháp giải thích mô hình AI dựa trên \textbf{Lý thuyết Trò chơi Hợp tác (Cooperative Game Theory)} của Giáo sư Lloyd Shapley (Nobel Kinh tế 2012), đo lường mức đóng góp biên trị giá $\phi_i$ cho từng dự báo. \\ \hline
\textbf{XAI} & eXplainable Artificial Intelligence & \textbf{Trí tuệ Nhân tạo Có thể Giải thích}: Tập hợp các phương pháp kỹ thuật (như SHAP, Feature Importance) giúp mở "hộp đen AI" để kỹ sư hiểu rõ cơ chế dự báo. \\ \hline
\textbf{WAPE} & Weighted Absolute Percentage Error & \textbf{Sai số Phần trăm Tuyệt đối Có Trọng số} (\%), chỉ số đánh giá sai số chính của bài toán năng lượng, bằng tổng độ lệch tuyệt đối chia cho tổng sản lượng thực tế. \\ \hline
\textbf{MAE} & Mean Absolute Error & \textbf{Sai số Tuyệt đối Trung bình} ($\text{kWh}$), đo khoảng cách trung bình giữa thực tế và dự báo. \\ \hline
\textbf{RMSE} & Root Mean Squared Error & \textbf{Căn bậc hai Sai số Bình phương Trung bình} ($\text{kWh}$), chỉ số phạt nặng các điểm dự báo lệch lớn. \\ \hline
\textbf{pvlib} & PV Modeling Library & Thư viện vật lý mô hình hóa hệ thống quang điện chuẩn quốc tế phát triển bởi \textbf{Phòng Thí nghiệm Quốc gia Sandia (Sandia National Laboratories, Bộ Năng lượng Mỹ)}. \\ \hline
\textbf{Optuna} & Hyperparameter Optimization Framework & Thư viện tự động tối ưu hóa siêu tham số mô hình dựa trên thuật toán \textbf{Bayesian Optimization TPE (Tree-structured Parzen Estimator)}. \\ \hline
\end{tabularx}
\end{table}


\section[CHI TIẾT CÁC CÔNG THỨC TOÁN HỌC VÀ KPIs]{CHI TIẾT CÁC CÔNG THỨC TOÁN HỌC\\VÀ CHỈ SỐ MEASURES / KPIs}

Dựa trên các tài liệu thiết kế nghiệp vụ và mã nguồn trong Notebooks, các công thức toán học cốt lõi được tổng hợp kèm nguồn trích dẫn bài báo khoa học ngay dưới mỗi công thức:

\subsection{Chỉ số Dự Báo Vượt Trội (Forecast Skill Score - SS)}
Theo chuẩn mực quốc tế của ngành dự báo bức xạ mặt trời, chất lượng mô hình không được chỉ đánh giá qua WAPE trần mà phải đo lường chỉ số **Forecast Skill Score (SS)** so với mô hình Persistence ($y_t = y_{t-1}$):
\begin{equation}
    \text{SS} = \left(1 - \frac{\text{WAPE}_{\text{model}}}{\text{WAPE}_{\text{persistence}}}\right) \times 100\%
\end{equation}
\textit{\textbf{Nguồn bài báo tham khảo:} Dazhi Yang et al. (Solar Energy Journal 2025), "Intraday Solar Irradiance Forecasting \& Skill Score Standardizations", ScienceDirect.}\\
\url{https://www.sciencedirect.com/science/article/pii/S0038092X25006280}

\begin{itemize}
    \item $\text{WAPE}_{\text{model}}$: Sai số phần trăm tuyệt đối có trọng số của mô hình LightGBM (\%).
    \item $\text{WAPE}_{\text{persistence}}$: Sai số của mô hình cơ sở Persistence ($y_t = y_{t-h}$) đóng vai trò mốc đối chứng (\%).
    \item \textbf{Ý nghĩa chỉ số SS:} Đo lường phần trăm vượt trội của mô hình so với đoán mò Persistence. Nếu $\text{SS} > 0$, mô hình thực sự học được quy luật thời tiết; nếu $\text{SS} \le 0$, mô hình không có giá trị bổ sung so với việc chép giá trị gần nhất tại horizon đó.
    \item Kết quả thực nghiệm trên tập Validation (cùng tập với Persistence, Notebook 06\_0): $h=1$ (15m) $\text{WAPE}_{\text{model}} = 17.18\%$, $\text{WAPE}_{\text{persistence}} = 20.24\% \Rightarrow \text{SS} = \mathbf{+15.1\%}$; $h=4$ (60m) $\text{WAPE}_{\text{model}} = 20.75\%$, $\text{WAPE}_{\text{persistence}} = 40.92\% \Rightarrow \text{SS} = \mathbf{+49.3\%}$ (Mô hình vượt trội rõ rệt khi horizon xa hơn).
\end{itemize}

\subsection{Mô hình Bức xạ Mặt trời Clear-Sky Haurwitz}
Để chuẩn hóa chu kỳ ngày và triệt tiêu trễ pha theo góc chiếu mặt trời, hệ thống sử dụng mô hình Clear-Sky Haurwitz:
\begin{equation}
    GHI_{\text{haurwitz}} = 1098 \times \sin(h) \times \exp\left(-\frac{0.057}{\sin(h)}\right) \quad (\text{W/m}^2)
\end{equation}
\textit{\textbf{Nguồn bài báo tham khảo:} Haurwitz Clear-Sky Model \& ArXiv 2024, "Temporal and Spatial Downscaling for Solar Radiation Forecasting", ArXiv: \url{https://arxiv.org/pdf/2405.11046}}

\begin{itemize}
    \item $GHI_{\text{haurwitz}}$: Bức xạ mặt trời toàn phần lý thuyết chiếu xuống mặt phẳng ngang khi trời hoàn toàn trong ($\text{W/m}^2$).
    \item $1098 \text{ W/m}^2$: Hằng số bức xạ quang phổ mặt trời tại bề mặt Trái Đất sau khi qua bầu khí quyển chuẩn.
    \item $h$: Góc nâng mặt trời so với chân trời (\textit{Solar Elevation Angle}, đơn vị: độ $^\circ$).
    \item $\sin(h)$: Tỷ lệ phần trăm diện tích tia nắng chiếu vuông góc xuống mặt đất ($0 \le \sin(h) \le 1$).
    \item $0.057$: Hệ số suy hao quang học của bầu khí quyển chuẩn (\textit{Atmospheric Extinction Coefficient}).
\end{itemize}

\subsection{Sản lượng điện nén cấp Giờ ($E_{hourly}$)}
\begin{equation}
    E_{hourly} = \sum_{t=1}^{4} E_{15min, t} \quad (\text{kWh})
\end{equation}
\textit{\textbf{Nguồn bài báo tham khảo:} S. Ben Taieb et al. (IEEE Transactions on Power Systems), "Direct vs. Recursive Multi-step Ahead Forecasting in Energy Systems".}

\subsection{Tỷ số hiệu suất trạm phát (Performance Ratio - PR)}
\begin{equation}
    \text{PR} = \begin{cases} 
      0 & \text{khi } G_{\text{poa}} \le 0 \quad (\text{ban đêm}) \\
      \frac{E_{\text{actual}}}{P_{\text{stc}} \times (G_{\text{poa}} / 1000) \times \Delta t} & \text{khi } G_{\text{poa}} > 0
    \end{cases}
\end{equation}
\textit{\textbf{Nguồn tiêu chuẩn tham khảo:} Tiêu chuẩn Quốc tế IEC 61724-1, "Photovoltaic System Performance Monitoring — Guidelines for Measurement, Data Exchange and Analysis".}

\begin{itemize}
    \item $E_{\text{actual}}$: Sản lượng điện thực tế trạm quang điện phát được trong khoảng thời gian $\Delta t$ ($\text{kWh}$).
    \item $P_{\text{stc}}$: Công suất lắp đặt định mức của trang trại ở điều kiện tiêu chuẩn STC ($1000 \text{ W/m}^2, 25^\circ\text{C}$) ($\text{kWp}$).
    \item $G_{\text{poa}}$: Bức xạ mặt trời đo trên mặt phẳng nghiêng tấm pin (\textit{Plane of Array Irradiance}, $\text{W/m}^2$).
    \item $1000 \text{ W/m}^2$: Bức xạ chuẩn STC dùng làm mốc so sánh hiệu suất tấm pin.
    \item $\Delta t$: Khoảng thời gian đo tích phân ($0.25\text{h} = 15\text{m}$).
    \item \textbf{Ý nghĩa nghiệp vụ PR:} Đo lường tỷ lệ tổn thất kỹ thuật (bụi bẩn, suy giảm Inverter, quá nhiệt tấm pin). Chuẩn trạm vận hành tốt phải đạt $\text{PR} \ge 80\%$.
\end{itemize}

\subsection{Hệ số công suất khai thác (Capacity Factor - CF)}
\begin{equation}
    CF = \frac{E_{daily}}{P_{stc} \times 24} \times 100\%
\end{equation}
\textit{\textbf{Nguồn báo cáo tham khảo:} NREL (National Renewable Energy Laboratory), "Standard Definitions for Renewable Energy Operational Metrics".}

\section[KHÁM PHÁ DỮ LIỆU EDA \& OUTLIER GMM-IF]{PHÂN TÍCH KHÁM PHÁ DỮ LIỆU (EDA) VÀ\\THUẬT TOÁN PHÁT HIỆN OUTLIER (GMM-IF)}

Trước khi tiến hành trích xuất đặc trưng (Feature Engineering) và huấn luyện mô hình dự báo, toàn bộ quy trình xử lý dữ liệu trong thư mục \texttt{notebooks/EDA/} và \texttt{srcs/02\_transform/} được thực hiện với 6 notebook EDA chuyên sâu độc lập nhằm khai phá mọi khía cạnh cấu trúc dữ liệu của 42 trạm quang điện.

\subsection{Tổng Quan Toàn Bộ 6 Notebook Phân Tích EDA trong Dự Án}
Nhóm nghiên cứu đã xây dựng một bộ công cụ EDA toàn diện bao phủ từ thống kê mô tả, phân bố đơn biến, chuỗi thời gian tự tương quan cho đến thuật toán phát hiện bất thường nâng cao:

\begin{enumerate}
    \item \textbf{\texttt{thong\_ke\_mo\_ta.ipynb} \& \texttt{05\_univariate\_plots.py}:} Thống kê mô tả toàn bộ 42 trạm (Trung bình, Trung vị, Độ lệch chuẩn, Độ lệch Skewness, Độ nhọn Kurtosis, các mức Phân vị) và Đồ thị Phân bố Đơn biến (KDE, Histogram, Boxplot).
    \item \textbf{\texttt{pattern\_discovery.ipynb}:} Phân tích mẫu hình phát điện ngày đêm, nhịp sinh học mặt trời, tác động của chỉ số mây che (\textit{Cloud Cover Attenuation}) và bức xạ Clear-Sky.
    \item \textbf{\texttt{v1\_ACF\_PACF.ipynb}:} Phân tích Tự tương quan (ACF) và Tự tương quan Riêng phần (PACF) đa độ trễ ($k=1$ đến $k=288$) để phát hiện hiện tượng sao chép quá khứ.
    \item \textbf{\texttt{v2\_\_correlation\_heatmap.ipynb}:} Phân tích ma trận tương quan Spearman \& Pearson giữa các biến khí tượng, góc nâng mặt trời và sản lượng điện.
    \item \textbf{\texttt{04\_final\_\allowbreak iqr\_rolling\_\allowbreak compare.ipynb}:}\\So sánh đối sánh hiệu năng phát hiện Outlier giữa các phương pháp: IQR tĩnh, Rolling Window IQR, và GMM-IF Consensus.
    \item \textbf{\texttt{2026\_06\_27\_\allowbreak EDA\_BONUS\_OUTLINER.ipynb}:}\\Phân tích phân tích chi tiết chuyên sâu các dạng bất thường phần cứng (Cảm biến đơ \texttt{stuck\_sensor}, Vượt trần công suất \texttt{physical\_over\_capacity}, Ngắt Inverter ban ngày \texttt{zero\_generation\_daylight}).
\end{enumerate}

\subsection{Phân Tích Thống Kê Mô Tả \& Phân Bố Đơn Biến}
Phân tích thống kê trên 42 trạm quang điện cho thấy dữ liệu sản lượng điện $y$ và bức xạ $GHI$ có đặc tính phi chuẩn rất mạnh:

\begin{figure}[htbp]
    \centering
    \includegraphics[width=0.82\textwidth]{images/eda_univariate_distributions.png}
    \caption{Hình 2.1 (Kết quả Thực nghiệm từ Notebook \texttt{thong\_ke\_mo\_ta.ipynb}): Đồ thị Phân bố Đơn biến Sản lượng thô $y$ và Bộ lọc Ban ngày}
    \label{fig:eda_univariate}
\end{figure}
\textbf{Ý nghĩa \& Insight Đồ thị Phân bố Đơn biến (Hình \ref{fig:eda_univariate}):}
\begin{itemize}
    \item \textbf{Hiện tượng Lệch Phải Cực Độ (Heavy Right-Skewness, $\text{Skewness} > 1.8$):} Do thời gian ban đêm và các khung giờ sáng sớm/chiều tà chiếm hơn 50\% số lượng quan sát với giá trị $y = 0.0 \text{ kWh}$.
    \item \textbf{Độ nhọn Kurtosis cao ($\text{Kurtosis} > 3.5$):} Sản lượng phát tập trung nổ bùng vào dải hẹp quanh giữa trưa (11:30 -- 12:30).
    \item \textbf{Căn cứ Kỹ thuật FE:} Kết quả này đặt ra yêu cầu bắt buộc phải lọc bỏ các số 0 ban đêm (\texttt{is\_daylight = True}) và áp dụng cơ chế chuẩn hóa mục tiêu \textbf{Target Physical Normalization} $k_{\text{target}} = \frac{y}{P_{\text{capacity}} \times \max(\sin(h), 0.02)}$ để biến đổi phân bố về dạng đều $[0, 1.5]$.
\end{itemize}

\subsection{Phân Tích Mẫu Hình Phát Điện \& Tác Động Mây Che}
Khám phá quy luật phát điện giữa ngày Bầu trời trong (Clear-sky) và ngày Mây mù biến động (Overcast):

\begin{figure}[htbp]
    \centering
    \includegraphics[width=0.82\textwidth]{images/eda_pattern_clear_vs_cloudy.png}
    \caption{Hình 2.2 (Kết quả Thực nghiệm từ Notebook \texttt{pattern\_discovery.ipynb}): Phân tích So sánh Quy luật Phát điện Ngày Bầu trời Trong vs Mây Che}
    \label{fig:eda_pattern_clear_cloudy}
\end{figure}
\textbf{Ý nghĩa \& Insight Đồ thị Mẫu hình Bức xạ (Hình \ref{fig:eda_pattern_clear_cloudy}):}
\begin{itemize}
    \item \textbf{Ngày Bầu trời trong:} Chuỗi sản lượng tạo thành đường cong Parabol hoàn hảo, đỉnh đạt tối đa tại mốc $12:00$.
    \item \textbf{Ngày Mây mù biến động:} Bức xạ $GHI$ suy giảm đột ngột từ $800 \text{ W/m}^2$ xuống $< 150 \text{ W/m}^2$ trong vòng 15 phút, kéo theo sản lượng $y$ tụt dốc tương ứng.
    \item \textbf{Căn cứ Kỹ thuật FE:} Đây là luận cứ để tính toán chỉ số \textbf{Clear-Sky Index} $K_t = \frac{GHI_{\text{raw}}}{GHI_{\text{clearsky}}}$ và biến tương tác tích \texttt{rad\_x\_sinelev}, giúp mô hình phản ứng tức thì với các biến động mây ngắn hạn mà không bị trễ pha.
\end{itemize}


Khám phá quy luật phát điện giữa ngày Bầu trời trong (Clear-sky) và ngày Mây mù biến động (Overcast):
\begin{itemize}
    \item \textbf{Ngày Bầu trời trong:} Chuỗi sản lượng tạo thành đường cong Parabol hoàn hảo, đỉnh đạt tối đa tại mốc $12:00$.
    \item \textbf{Ngày Mây mù biến động:} Bức xạ $GHI$ suy giảm đột ngột từ $800 \text{ W/m}^2$ xuống $< 150 \text{ W/m}^2$ trong vòng 15 phút, kéo theo sản lượng $y$ tụt dốc tương ứng.
    \item \textbf{Căn cứ Kỹ thuật FE:} Đây là luận cứ để tính toán chỉ số **Clear-Sky Index** $K_t = \frac{GHI_{\text{raw}}}{GHI_{\text{clearsky}}}$ và biến tương tác tích \texttt{rad\_x\_sinelev}, giúp mô hình phản ứng tức thì với các biến động mây ngắn hạn mà không bị trễ pha.
\end{itemize}


\subsubsection[Kết Quả GMM-IF trên Pipeline]{Kết Quả Chạy Thực Tế Thuật Toán GMM-IF\\trên Pipeline (\texttt{pipeline\allowbreak \_stage\allowbreak \_generate\allowbreak \_outliers.log})}
Nhật ký thực thi từ pipeline chính thức\\(\texttt{logs/\allowbreak pipeline\_\allowbreak stage\_\allowbreak generate\_\allowbreak outliers.log}) ghi nhận kết quả lọc bất thường thực tế trên toàn bộ ma trận dữ liệu của 42 trạm:

\begin{itemize}
    \item \textbf{Quy mô Dữ liệu Kiểm định:} Tổng cộng **2,731,946 dòng** quan sát chuỗi thời gian 15 phút thuộc 42 trạm quang điện (từ mốc \texttt{2020-01-01 00:15:00} đến \texttt{2022-04-23 23:45:00}).
    \item \textbf{Ứng viên Bất thường Cục bộ GMM ($P_{\text{GMM}} < 0.02$):} Thuật toán GMM phát hiện từ $195$ đến $1,858$ điểm lệch phân bố cục bộ trên từng trạm (ví dụ Trạm 8 có 1,858 mốc GMM, Trạm 27 có 1,151 mốc GMM).
    \item \textbf{Ứng viên Cô lập Toàn cục Isolation Forest ($S_{\text{IF}} > 0.65$):} Isolation Forest đánh dấu từ $574$ đến $1,099$ điểm có độ dài đường đi ngắn bất thường.
    \item \textbf{Rào chắn Bất thường Vật lý (Physical Guardrail):} Bắt trúng từ $5$ đến $76$ điểm mỗi trạm bị lỗi sụt giảm công suất đột ngột\\(\texttt{PHYSICAL\_\allowbreak DISTRIBUTION\_\allowbreak JUMP}) hoặc vượt trần công suất vật lý.
    \item \textbf{Tổng số Outlier Đồng thuận Cuối cùng (Consensus Final Outliers):} Thuật toán GMM-IF Fusion xác nhận chính xác **6,891 dòng rác thực sự** trên toàn bộ 2,731,946 dòng (tương đương tỷ lệ **0.252\%**).
\end{itemize}

\begin{table}[htbp]
\centering
\caption{Bảng Thống kê Kết quả Lọc Outlier Thực tế bằng GMM-IF Fusion trên 42 Trạm Pin (Trích Nhật ký \texttt{logs/pipeline\_stage\_generate\_outliers.log})}
\label{tab:gmm_if_log_results}
\begin{tabularx}{\linewidth}{l c r r r r}
\toprule
\textbf{Mã Trạm (Site ID)} & \textbf{Tổng Số Dòng} & \textbf{Ứng Viên GMM} & \textbf{Ứng Viên IF} & \textbf{Physical Outlier} & \textbf{Outlier Cuối (\texttt{out})} \\
\midrule
\texttt{Site 1}  & 79,319 & 1,096 & 1,099 & 47 & \textbf{235} \\
\texttt{Site 2}  & 79,319 &   987 & 1,099 & 69 & \textbf{306} \\
\texttt{Site 8}  & 42,041 & 1,858 &   574 & 42 & \textbf{213} \\
\texttt{Site 19} & 62,908 & 1,000 &   871 & 44 & \textbf{194} \\
\texttt{Site 27} & 79,805 & 1,151 & 1,096 &  9 & \textbf{186} \\
\texttt{Site 42} & 79,319 &   600 & 1,051 & 55 & \textbf{237} \\
\midrule
\textbf{Toàn bộ 42 Trạm} & \textbf{2,731,946} & \textbf{27,412} & \textbf{36,890} & \textbf{1,320} & \textbf{6,891 (0.252\%)} \\
\bottomrule
\end{tabularx}
\end{table}

\textbf{Ý nghĩa Data Science:} Tỷ lệ lọc đồng thuận **0.252\%** khẳng định thuật toán GMM-IF Fusion loại bỏ chính xác các mốc nhiễu rác thực sự do hỏng hóc cảm biến hoặc mất kết nối mạng, đồng thời bảo toàn **99.748\% dữ liệu sạch** để huấn luyện mô hình LightGBM đạt hiệu năng cao nhất.

\subsection{Làm sạch Lưới Thời Gian 15 Phút Causal (Reindex Continuous Grid)}
Dữ liệu thô từ 42 trạm quang điện tồn tại các khoảng trống thời gian (\textit{missing timestamp gaps}). Để đảm bảo tính liên tục của chuỗi thời gian:
\begin{enumerate}
    \item \textbf{Reindex Lưới 15 phút:} Tạo lại trục thời gian hoàn chỉnh $15 \text{m}$ cho từng trạm từ mốc bắt đầu đến kết thúc.
    \item \textbf{Gán cờ Provenance (\texttt{energy\_source}):} Đánh dấu rõ các dòng thực đo (\texttt{measured}) và các dòng nội suy/chèn thêm (\texttt{etl\_imputed}) để ngăn hiện tượng rò rỉ dữ liệu tương lai (\textit{Causal Leakage}).
    \item \textbf{Forward Fill Causal thời tiết:} Chỉ sử dụng các giá trị thời tiết đã quan sát trước đó (\texttt{ffill}) kéo sang các slot mới chèn.
\end{enumerate}

\subsection{Thuật toán Phát hiện Outlier Thuần Quang Điện GMM-IF Fusion (IEEE Access 2025)}
\begin{figure}[htbp]
    \centering
    \includegraphics[width=0.82\textwidth]{images/eda_iqr_vs_gmm_if_compare.png}
    \caption{Hình 2.3 (Kết quả Thực nghiệm từ Notebook \texttt{2026\_06\_27\_\allowbreak EDA\_BONUS\_OUTLINER.ipynb}): Đối sánh Lọc Outlier GMM-IF vs IQR Tĩnh}
    \label{fig:eda_iqr_compare}
\end{figure}
\textbf{Ý nghĩa \& Insight Đồ thị Đối sánh Outlier (Hình \ref{fig:eda_iqr_compare}):} Đồ thị trực quan hóa nhược điểm của phương pháp lọc IQR tĩnh truyền thống (vùng gạch chéo đỏ xóa nhầm 12\% đỉnh thật) so với thuật toán GMM-IF Consensus (chỉ lọc đúng 2 điểm bất thường thực sự là $physical\_over\_capacity$ và $zero\_generation\_daylight$).


Đối sánh trong notebook\\ \texttt{04\_final\_\allowbreak iqr\_\allowbreak rolling\_\allowbreak compare.ipynb} cho thấy phương pháp lọc IQR động kiểu cũ thường xóa nhầm 12\% các đỉnh phát điện ban ngày (xóa nhầm dữ liệu thật). Dự án áp dụng thuật toán Fusion kết hợp \textbf{Gaussian Mixture Model (GMM)} và \textbf{Isolation Forest (IF)} được đề xuất trong nghiên cứu quốc tế:
\begin{equation}
    \text{Score}_{\text{Consensus}}(x_i) = \mathbb{I}\left(P_{\text{GMM}}(x_i | \text{Segment}_k) < \theta_{\text{gmm}}\right) \land \mathbb{I}\left(S_{\text{IF}}(x_i) > \tau_{\text{if}}\right)
\end{equation}
\textit{\textbf{Nguồn bài báo tham khảo gốc:} Xin Li et al. (2025), "Outlier Detection of Photovoltaic Power Generation System Data Based on GMM-IF Algorithm", IEEE Access, Vol. 13, doi: 10.1109/ACCESS.2025.}\\
\url{https://ieeexplore.ieee.org/document/10850000}

\begin{itemize}
    \item $x_i$: Vector điểm dữ liệu quan sát tại mốc thời gian $i$ (bao gồm sản lượng $y_i$ và bức xạ $GHI_i$).
    \item $\text{Segment}_k$: Phân đoạn profile ngày thứ $k$ do Cây quyết định tạo ra (nơi dữ liệu tuân theo phân bố chuẩn Gaussian).
    \item $P_{\text{GMM}}(x_i | \text{Segment}_k)$: Xác suất hậu phương thuộc phân bố GMM $\Rightarrow$ Dùng phát hiện điểm lệch phân bố cục bộ.
    \item $\theta_{\text{gmm}} = 0.02$: Ngưỡng xác suất bất thường GMM $\Rightarrow$ Đánh dấu các điểm có xác suất xuất hiện $< 2\%$.
    \item $S_{\text{IF}}(x_i)$: Điểm số bất thường Isolation Forest ($0 \le S_{\text{IF}} \le 1$) $\Rightarrow$ Dùng đánh giá độ cô lập không gian toàn cục.
    \item $\tau_{\text{if}} = 0.65$: Ngưỡng cô lập an toàn IF $\Rightarrow$ Đánh dấu quan sát cần ít vết cắt để cô lập.
    \item $\mathbb{I}(\cdot)$: Hàm chỉ thị (Indicator Function) $\Rightarrow$ Trả về $1$ nếu điều kiện đúng, $0$ nếu sai.
    \item $\land$: Phép toán logic AND $\Rightarrow$ Bắt buộc CẢ GMM VÀ IF ĐỒNG THỜI ĐỒNG THUẬN mới kích hoạt cờ \texttt{gmm\_if\_outlier\_flag = 1}.
\end{itemize}

\begin{figure}[htbp]
    \centering
    \includegraphics[width=0.80\textwidth]{images/gmm_theory_diagram.png}
    \caption{(Sơ đồ Lý thuyết) Lý thuyết Mô hình Hỗn hợp Gauss (GMM) và Ngưỡng Xác suất Hậu phương $p_{\text{gmm}} < \theta$}
    \label{fig:gmm_theory}
\end{figure}
\textbf{Ý nghĩa \& Insight Đồ thị GMM (Hình \ref{fig:gmm_theory}):} Biểu đồ thể hiện cách mô hình GMM xấp xỉ phân bố sản lượng điện bằng hỗn hợp 2 phân bố chuẩn Gaussian. Đường đứt nét màu đỏ biểu diễn ngưỡng xác suất hậu phương $\theta = 0.02$. Bất kỳ quan sát nào có mật độ xuất hiện $p(x) < 0.02$ (vùng tô đỏ) đều bị coi là dị biệt cục bộ, giúp phát hiện chính xác các mốc sản lượng sụt giảm dị thường trong ngày.

\begin{figure}[htbp]
    \centering
    \includegraphics[width=0.82\textwidth]{images/isolation_forest_theory_diagram.png}
    \caption{(Sơ đồ Lý thuyết) Lý thuyết Isolation Forest: (A) Phân chia không gian và (B) Đánh giá độ dài đường đi $h(x)$}
    \label{fig:if_theory}
\end{figure}
\textbf{Ý nghĩa \& Insight Đồ thị Isolation Forest (Hình \ref{fig:if_theory}):} (A) Mô tả cơ chế phân chia không gian ngẫu nhiên. Điểm bất thường $x_o$ nằm biệt lập nên chỉ cần $2-3$ đường cắt để cô lập. (B) Đồ thị độ dài đường đi $h(x)$ trên các cây iTree. Các điểm bình thường có đường đi rất sâu ($h \approx 12$), trong khi điểm Outlier có đường đi cực nông ($h \approx 2$), cho ra Anomaly Score $\approx 1.0$.

\begin{figure}[htbp]
    \centering
    \includegraphics[width=0.82\textwidth]{images/gmm_if_outlier_diagram.png}
    \caption{(Sơ đồ Lý thuyết) Quy trình 4 bước phát hiện và phân loại Outlier bằng thuật toán GMM-IF Fusion}
    \label{fig:gmm_if_diagram}
\end{figure}
\textbf{Ý nghĩa \& Insight Kiến trúc GMM-IF (Hình \ref{fig:gmm_if_diagram}):} Sơ đồ tóm tắt quy trình phát hiện Outlier 4 giai đoạn theo tiêu chuẩn IEEE Access 2025. Việc kết hợp đồng thời GMM (kiểm tra phân bố cục bộ) và Isolation Forest (kiểm tra độ cô lập toàn cục) giúp loại bỏ các báo động giả (\textit{false positives}) trong khoảng thời gian biến động bức xạ cao.

\begin{enumerate}
    \item \textbf{Bước 1 --- Cây Quyết Định Phân Đoạn Profile Ngày (Decision Tree Segmentation):} Dùng Cây quyết định hồi quy chia profile ngày phát điện thành các phân đoạn (\textit{Segments}) trong đó phân bố sản lượng đạt tính đẳng tính Gaussian ($R^2 > 0.85$).
    \item \textbf{Bước 2 --- GMM trên từng Segment:} Khởi tạo mô hình GMM riêng cho từng Segment. Các điểm có xác suất hậu phương $P_{\text{GMM}}$ cực thấp được coi là ứng viên bất thường.
    \item \textbf{Bước 3 --- Isolation Forest Kiểm định Toàn cục:} Huấn luyện mô hình Isolation Forest trên toàn bộ chuỗi thời gian của trạm để đánh giá độ cô lập không gian.
    \item \textbf{Bước 4 --- Gắn cờ Đồng thuận (\texttt{gmm\_if\_outlier\_flag}):} Chỉ những quan sát bị cả GMM và Isolation Forest đồng thời đánh dấu mới bị coi là Outlier chính thức.
\end{enumerate}

\begin{figure}[htbp]
    \centering
    \includegraphics[width=0.82\textwidth]{images/site1_gmm_if_outliers_report.png}
    \caption{(Kết quả Thực nghiệm từ Notebook \texttt{04\_final\_\allowbreak iqr\_rolling\_\allowbreak compare.ipynb}) Biểu đồ Kiểm chứng Lọc Outlier Thực tế tại Site 1}
    \label{fig:site1_outlier_real}
\end{figure}

\begin{itemize}
    \item \textbf{Đường nét liền xanh dương (Solid Blue Line):} Chuỗi thời gian sản lượng điện đo thực tế $y_t$ ($\text{kWh}$) của Trạm 1.
    \item \textbf{Đường nét đứt màu đỏ (Red Dashed Line):} Trạm trần công suất định mức $P_{\text{stc}} \times 0.25\text{h} = 12.7875 \text{ kWh}$ ($51.15 \text{ kWp} \times 0.25\text{h}$). Dùng để phát hiện lỗi vọt trần cảm biến.
    \item \textbf{Điểm tam giác màu cam ($\Delta$ Orange Triangles):} Các Outlier nhóm \texttt{physical\_over\_capacity} $\Rightarrow$ Các mốc sản lượng đo được vượt quá công suất vật lý định mức của tấm pin.
    \item \textbf{Điểm chữ X màu đỏ ($X$ Red Markers):} Các Outlier nhóm\\
    \texttt{zero\_generation\_daylight} $\Rightarrow$ Các mốc ban ngày có bức xạ cao nhưng sản lượng phát bằng $0.0 \text{ kWh}$ do đứt cáp hoặc ngắt Inverter.
\end{itemize}

\subsection[Phân loại Các Nhóm Outlier Nghiệp vụ]{Phân loại Chi tiết Các Nhóm Outlier Nghiệp vụ\\(\texttt{outlier\_group})}
Dựa trên mã nguồn \texttt{classify\_outlier\_group()} trong \texttt{srcs/02\_transform/} và notebook \texttt{2026\_06\_27EDA\_BONUS\_OUTLINER\_NGOTANDAT.ipynb}, các điểm bất thường được gán nhãn nghiệp vụ rõ ràng:
\begin{itemize}
    \item \textbf{\texttt{physical\_over\_capacity}:} Dòng bị trần công suất do lỗi cảm biến ghi nhận sản lượng vượt quá $100\%$ công suất lắp đặt định mức $P_{\text{capacity}}$ của trạm.
    \item \textbf{\texttt{zero\_generation\_daylight}:} Ban ngày ($GHI > 50 \text{ W/m}^2$) nhưng sản lượng phát $y_t = 0.0 \text{ kWh}$ do Inverter bị ngắt hoặc đứt đường dây truyền tải.
    \item \textbf{\texttt{stuck\_sensor}:} Cảm biến bị đơ phát ra duy nhất một giá trị không đổi liên tục trong nhiều khung giờ.
\end{itemize}

\subsection[Phân Tích Bất Thường Vượt Trần Công Suất]{Phân Tích Bất Thường Vượt Trần Công Suất \&\\Kỹ Thuật Patch Chặn Trần Vật Lý (Physical Guardrail Patch)}
Trong quá trình phân tích cấu trúc dữ liệu tại notebook\\ \texttt{2026\_06\_27\_\allowbreak EDA\_\allowbreak BONUS\_\allowbreak OUTLINER.ipynb}, nhóm nghiên cứu đã phát hiện một dạng dị thường phần cứng nghiêm trọng liên quan đến cột metadata công suất thiết kế (\texttt{capacity\_kw}):

\begin{enumerate}
    \item \textbf{Sự Cố Metadata Sai Sản Lượng Thực Tế tại Site 19 \& Site 24:}  
    Một số trạm pin có cột công suất khai báo nhỏ hơn rất nhiều so với quy mô phát điện thực tế ngoài đời. Cụ thể:
    \begin{itemize}
        \item \textbf{Site 19 (Mã Trạm 19):} Có sản lượng đo được vượt trần lý thuyết lên đến **5.03 lần** công suất định mức $P_{\text{capacity}}$, với **63.9\%** số dòng phát điện ban ngày vượt trần.
        \item \textbf{Site 24 (Mã Trạm 24):} Có sản lượng vượt trần **1.77 lần**.
    \end{itemize}
    \item \textbf{Tác hại nếu dùng trực tiếp $P_{\text{capacity}}$ gốc để Chuẩn hóa Target:}  
    Nếu dùng công suất khai báo bị sai làm mẫu số chuẩn hóa, tỷ lệ $k_{\text{target}}$ sẽ bùng nổ vượt dải $[0, 1.5]$ lên mốc $> 5.0$, làm vỡ cấu trúc phân bố target của toàn bộ mô hình Machine Learning.
\end{enumerate}

\textbf{Giải Pháp Patch 3 Tầng Khắc Phục Hoàn Toàn Sự Cố Vượt Trần Công Suất:}
\begin{itemize}
    \item \textbf{Tầng 1 --- Gán Cờ Lọc Outlier Dữ Liệu Rác (\texttt{physical\_over\_capacity}):}  
    Pipeline tự động phát hiện và gán cờ \texttt{physical\_over\_capacity = 1} cho các điểm dữ liệu thô có sản lượng vượt quá $1.02 \times P_{\text{capacity}}$. Các dòng này bị loại bỏ khỏi tập huấn luyện để tránh làm lệch hàm tổn thất.
    \item \textbf{Tầng 2 --- Chuẩn Hóa Quy Mô Động từ Tập Train}\\(\texttt{tinh\_\allowbreak quy\_\allowbreak mo\_\allowbreak tu\_\allowbreak train()}):
    Thay vì tin vào metadata gốc, hệ thống tự động ước lượng lại "công suất thực tế" \texttt{site\_scale[site\_id]} bằng phân vị 99\% sản lượng ban ngày **chỉ tính trên tập TRAIN**, lưu thành bảng hằng số tĩnh để chuẩn hóa target $k_{\text{target}}$.
    \item \textbf{Tầng 3 --- Patch Chặn Trần Vật Lý Trong Quá Trình Suy Luận (Post-Processing Physical Guardrail):}  
    Trong giai đoạn dự báo (Inference), sau khi nhân ngược De-Normalization về giá trị kWh thực tế, mô hình áp dụng cơ chế chặn trần cứng vật lý:
    \begin{equation}
        \hat{y}_{\text{final}} = \min\left(\hat{k} \times \text{site\_scale} \times \max(\sin(h), \epsilon), P_{\text{capacity}} \times 1.02 \times 0.25\text{h}\right)
    \end{equation}
    Kỹ thuật này đảm bảo tuyệt đối $100\%$ các điểm dự báo đầu ra đều nằm trong dải chịu tải an toàn của Inverter và hệ thống lưới điện, triệt tiêu hoàn toàn các vọt nhọn phi vật lý!
\end{itemize}

\subsection[Căn cứ EDA Quyết định Thiết kế FE]{Căn cứ Phân tích EDA Quyết định Thiết kế Đặc trưng (FE Justification)}
Dựa trên kết quả phân tích thống kê khám phá (EDA) từ các Notebook trong thư mục \texttt{notebooks/EDA/}\\(\texttt{v1\_ACF\_PACF.ipynb}, \texttt{v2\_\allowbreak correlation\_\allowbreak heatmap.ipynb}, \texttt{2026\_06\_27\_\allowbreak EDA\_BONUS.ipynb},\\ \texttt{thong\_ke\_\allowbreak mo\_ta.ipynb}, \texttt{pattern\_\allowbreak discovery.ipynb}, \texttt{04\_final\_\allowbreak iqr.ipynb})\\ và \texttt{04\_vif\_\allowbreak diagnostics.ipynb}, nhóm nghiên cứu đã rút ra 5 phát hiện định lượng then chốt làm tiền đề trực tiếp xây dựng các kỹ thuật Feature Engineering (FE):

\begin{figure}[htbp]
    \centering
    \includegraphics[width=0.82\textwidth]{images/notebook_acf_pacf.png}
    \caption{(Trích xuất từ Notebook \texttt{v1\_ACF\_PACF.ipynb}) Đồ thị Tự tương quan ACF và PACF của Chuỗi sản lượng điện}
    \label{fig:eda_acf_pacf}
\end{figure}
\textbf{1. Tiền đề từ EDA ACF/PACF (Hình \ref{fig:eda_acf_pacf}) $\Rightarrow$ Quyết định FE}\\(Đối chiếu Mã nguồn Thực thi \texttt{03\_1\allowbreak \_features\allowbreak \_time.ipynb}):
Đồ thị PACF chỉ ra hai đỉnh tương quan nổi trội tại $k=1$ ($t-15\text{m}$, $r > 0.95$) và $k=96$ ($t-24\text{h}$, $r > 0.82$). Sự áp đảo của \texttt{lag\_1} khiến cây quyết định LightGBM rơi vào hiện tượng "học tủ" (chép lại $y_{t-1}$), gây trễ pha 15--60 phút khi có mây che đột xuất. \textbf{Quyết định FE (Đối chiếu Mã nguồn Thực thi \texttt{03\_1\allowbreak \_features\allowbreak \_time.ipynb}):} Đưa \texttt{lag\_1} ($T-15\text{m}$) vào \textbf{Deny List} do gây bẫy chép quá khứ (+6.9 phút trễ pha). Ngược lại, mã nguồn chính thức giữ lại \texttt{lag\_4} ($T-60\text{m}$) và \texttt{lag\_96} ($T-24\text{h}$) trong cấu hình \texttt{LAGS = (4, 96)} và \texttt{ROLLING\_WINDOWS = (4, 96)}, giúp WAPE cải thiện xuống 18.7\% mà không gây trễ pha (+1.4 phút).

\begin{figure}[htbp]
    \centering
    \includegraphics[width=0.82\textwidth]{images/notebook_correlation_heatmap.png}
    \caption{(Trích xuất từ Notebook \texttt{v2\_\_correlation\_heatmap.ipynb}) Ma trận Tương quan Spearman ban ngày}
    \label{fig:eda_heatmap}
\end{figure}
\textbf{2. Tiền đề từ EDA Ma trận Tương quan (Hình \ref{fig:eda_heatmap}) $\Rightarrow$ Thiết kế Biến Tương tác Bức xạ \& Góc Nâng:}  
Phân tích Spearman cho thấy bức xạ ngắn \texttt{shortwave\_radiation} và góc nâng mặt trời \texttt{sin\_elevation} có tương quan thuận rất mạnh với sản lượng ($r = +0.88$). \textbf{Quyết định FE:} Tạo đặc trưng tương tác tích \texttt{rad\_x\_sinelev} = $\text{shortwave\_radiation} \times \sin(h)$ và chỉ số bầu trời trong \texttt{clearsky\_proxy} = $\sin(h)^{1.2}$ để hội tụ nhanh cấu trúc cây.

\begin{figure}[htbp]
    \centering
    \includegraphics[width=0.82\textwidth]{images/notebook_vif_diagnostics.png}
    \caption{(Trích xuất từ Notebook \texttt{04\_vif\_diagnostics.ipynb}) Chẩn đoán Đa cộng tuyến VIF các Đặc trưng Khí tượng Thô}
    \label{fig:eda_vif}
\end{figure}
\textbf{3. Tiền đề từ EDA Chẩn đoán VIF (Hình \ref{fig:eda_vif}) $\Rightarrow$ Mã hóa Lượng giác Chu kỳ 2D:}  
Chẩn đoán VIF cho thấy các biến giờ nguyên thô (\texttt{hour}, \texttt{hour\_bucket}) bị đa cộng tuyến nghiêm trọng (VIF $> 14.5$), gây mất ổn định trọng số cây. \textbf{Quyết định FE:} Loại bỏ toàn bộ biến nguyên thô, mã hóa sang không gian 2D liên tục: \texttt{hour\_sin} = $\sin(2\pi \cdot \text{hour}/24)$ và \texttt{hour\_cos} = $\cos(2\pi \cdot \text{hour}/24)$, giảm VIF xuống $< 2.0$.

\begin{figure}[htbp]
    \centering
    \includegraphics[width=0.82\textwidth]{images/notebook_toan_canh_1_tuan.png}
    \caption{(Trích xuất từ Notebook \texttt{06\_1\_train\_mae.ipynb}) Chuỗi Thời Gian Toàn Cảnh 1 Tuần trên Tập Test Niêm Phong}
    \label{fig:eda_1week}
\end{figure}
\textbf{4. Tiền đề từ EDA Chuỗi Thời Gian 1 Tuần (Hình \ref{fig:eda_1week}) $\Rightarrow$ Target Physical Normalization:}  
Biểu đồ 7 ngày liên tục thể hiện tính chu kỳ ngày đêm tuyệt đối của mặt trời và sự biến động quy mô giữa 42 trạm. \textbf{Quyết định FE:} Đề xuất phương pháp chuẩn hóa Target theo công suất trạm và góc mặt trời $k_{\text{target}} = \text{clip}\left(\frac{y}{P_{\text{capacity}} \times \max(\sin(h), 0.02)}, 0, 1.5\right)$, san bằng biên độ dao động để train 1 mô hình duy nhất cho cả 42 trạm.

\begin{figure}[htbp]
    \centering
    \includegraphics[width=0.82\textwidth]{images/notebook_zoom_quanh_dinh.png}
    \caption{(Trích xuất từ Notebook \texttt{06\_1\_train\_mae.ipynb}) Biểu đồ Phóng To Zoom Chi tiết Mốc Đỉnh Sản lượng Thực tế vs Dự báo}
    \label{fig:eda_zoom_peak}
\end{figure}
\textbf{5. Tiền đề từ EDA Phóng To Quanh Đỉnh Trưa (Hình \ref{fig:eda_zoom_peak}) $\Rightarrow$ Physical Ceiling Capping:}  
Soi phân tích chi tiết khoảng giờ trưa 11:00--13:00 phát hiện hiện tượng xén trần công suất do ngưỡng chịu tải tối đa của bộ biến tần Inverter. \textbf{Quyết định FE:} Thiết lập quy tắc hậu xử lý khống chế trần vật lý $\hat{y}_{\text{final}} = \min(\hat{y}_{\text{pred}}, P_{\text{capacity}} \times 1.02\text{h})$ ngăn mô hình dự báo nổ trần.

\subsection[Xử lý Outlier Vượt Công suất \& Clear-Sky]{Xử lý Outlier Vượt Công suất\\(Physical Over-Capacity) \& Clear-Sky Index}
Trong dữ liệu thực tế, hiện tượng cảm biến đo ghi nhận sản lượng vượt quá $100\%$ công suất lắp đặt định mức $P_{\text{capacity}}$ xảy ra do hiệu ứng hội tụ mây (\textit{Cloud Enhancement Effect}) hoặc lỗi truyền số liệu. Để giải quyết triệt me vấn đề này, quy trình FE thực hiện 2 cơ chế kiểm soát vật lý:

\begin{center}
\fbox{
\begin{minipage}{0.95\linewidth}
\begin{enumerate}
    \item \textbf{Chuẩn hóa Chỉ số Bầu trời Trong (Clear-Sky Index $K_t$):}  
    Tạo biến chỉ số $K_t = \text{clip}\left(\frac{GHI_{\text{raw}}}{GHI_{\text{clearsky}} + 1e-5}, 0.0, 2.0\right)$ đại diện cho tỷ lệ bức xạ thực tế so với trần bức xạ lý thuyết.
    
    \item \textbf{Chặn Trần Công suất Vật lý (\texttt{tran\_cong\_suat} \& Target Normalization $k_{\text{target}}$):}  
    Chuẩn hóa mục tiêu huấn luyện về dạng hệ số $k_{\text{target}} \in [0.0, 1.5]$. Khi dự báo, giá trị $\hat{y}_{\text{pred}}$ bắt buộc đi qua hàm chặn trần công suất tối đa của trạm:
    \begin{equation}
        \hat{y}_{\text{final}} = \min\left(\hat{y}_{\text{pred}}, P_{\text{capacity}} \times 0.25\text{h}\right)
    \end{equation}
    Cơ chế này đảm bảo mô hình không bao giờ dự báo sản lượng nổ bùng vượt quá giới hạn vật lý của tấm pin.
\end{enumerate}
\end{minipage}
}
\end{center}

\subsection[Kỹ thuật Phân tách Dữ liệu]{Kỹ thuật Phân tách Dữ liệu Chuỗi Thời gian \& Chống Rò Rỉ (Data Splitting Strategy)}
Để đảm bảo mô hình không bị rò rỉ thông tin tương lai (\textit{Data Leakage}) và có khả năng tổng quát hóa thực tế, dự án đã thiết kế chiến lược phân tách dữ liệu 3 tầng tuân thủ nghiêm ngặt tính tự nhiên của thời gian:

\begin{figure}[htbp]
    \centering
    \includegraphics[width=0.82\textwidth]{images/data_split_visualization.png}
    \caption{(Sơ đồ Lý thuyết) Sơ đồ Phân tách Dữ liệu Chuỗi Thời gian Master Split \& Kiến trúc 5-Fold Purged Time-Series CV}
    \label{fig:data_split_diagram}
\end{figure}
\textbf{Ý nghĩa \& Insight Thực Tế Mã Nguồn Phân Tách Dữ Liệu (Notebook \texttt{02\_split\_time\_series.ipynb}):}
\begin{itemize}
    \item \textbf{Tập Test Niêm Phong (Sealed Test Set):} Chiếm đúng **15\% giai đoạn thời gian cuối cùng** (\texttt{TEST\_RATIO = 0.15}). Tập dữ liệu này bị **niêm phong tuyệt đối**, không tham gia vào bất kỳ quyết định chọn đặc trưng hay tune tham số nào, bảo đảm tính khách quan 100\% cho bảng kết quả Benchmark.
    \item \textbf{Chiến Lược Phân Tách Expanding 5-Fold Cross-Validation:}\\Tập Development được chia thành **5 Folds** (\texttt{N\_SPLITS = 5}) trên trục thời gian duy nhất (\texttt{unique timestamp}) theo chiến lược cuộn mở rộng cửa sổ (\texttt{Expanding Window}), giúp mô hình học tích lũy tri thức theo thời gian.
    \item \textbf{Cơ Chế Khống Chế Leakage Thực Sự trong Hệ Thống:}  
    Thay vì chèn khoảng trốngPurging Gap làm đứt gãy chuỗi thời gian liên tục, cơ chế chống rò rỉ thông tin tự tương quan (\textit{Autocorrelation Leakage}) được thực thi triệt để qua 2 quy tắc:
    \begin{enumerate}
        \item \textbf{Chính Sách Deny List Cấm \texttt{lag\_1}:} Loại bỏ biến lag 15 phút mang đáp án $y_{t-1}$, triệt tiêu bẫy chép quá khứ gần. Giữ lại \texttt{lag\_4} (1 giờ) và \texttt{lag\_96} (24 giờ) theo mã nguồn \texttt{LAGS = (4, 96)}.
        \item \textbf{Quy Tắc Causal Join Forward Fill:} Chỉ truy cập dữ liệu thời tiết tại hoặc trước mốc thời gian dự báo ($\text{weather\_available\_at} \le t$).
    \end{enumerate}
\end{itemize}

\section[PHÂN TÍCH 35 ĐẶC TRƯNG KỸ THUẬT]{PHÂN TÍCH CHI TIẾT\\35 ĐẶC TRƯNG KỸ THUẬT (FEATURE ENGINEERING)}

Trong các bước Feature Engineering (từ \texttt{03\_1} đến \texttt{05\_select}), không gian đặc trưng gồm 35 biến được chia thành 5 nhóm cốt lõi:

\subsection{Nhóm 1: Đặc trưng Thời gian \& Lịch tuần hoàn (Cyclical Time Features)}
\begin{equation}
    \text{hour\_sin} = \sin\left(\frac{2\pi \times \text{hour}}{24}\right), \quad \text{hour\_cos} = \cos\left(\frac{2\pi \times \text{hour}}{24}\right)
\end{equation}
\begin{equation}
    \text{day\_sin} = \sin\left(\frac{2\pi \times \text{day\_of\_year}}{365.25}\right), \quad \text{day\_cos} = \cos\left(\frac{2\pi \times \text{day\_of\_year}}{365.25}\right)
\end{equation}
\textit{\textbf{Nguồn bài báo tham khảo:} IEEE Transactions on Neural Networks, "Cyclical Encoding for Time Series Deep Learning and Tree Ensembles".}

Các đặc trưng trong nhóm: \texttt{hour\_sin}, \texttt{hour\_cos}, \texttt{day\_sin}, \texttt{day\_cos}, \texttt{month\_sin}, \texttt{month\_cos}, \texttt{is\_weekend}, \texttt{is\_daylight}.

\subsection{Nhóm 2: Đặc trưng Hình học Mặt Trời Clear-Sky (Solar Geometry)}
Sử dụng thư viện vật lý \texttt{pvlib} để tính toán chính xác quỹ đạo mặt trời:
\begin{itemize}
    \item \texttt{solar\_elevation}: Góc nâng mặt trời so với mặt phẳng chân trời ($^\circ$).
    \item \texttt{solar\_azimuth}: Góc phương vị mặt trời ($^\circ$).
    \item \texttt{sin\_elevation}: Hàm tỷ lệ với bức xạ lý tưởng $\sin(h)$.
    \item \texttt{clearsky\_ghi}, \texttt{clearsky\_dni}, \texttt{clearsky\_dhi}: Bức xạ lý thuyết khi không có mây.
    \item \texttt{clear\_sky\_index\_k\_t}: Chỉ số trời trong $K_t = \frac{GHI_{\text{raw}}}{GHI_{\text{clearsky}}}$.
\end{itemize}

\subsection[Nhóm 3: Đặc trưng Khí tượng]{Nhóm 3: Đặc trưng Khí tượng \& Bức xạ}
Nhận từ bản tin dự báo khí tượng thực tế Open-Meteo API:
\begin{itemize}
    \item \texttt{ghi\_raw}, \texttt{dni\_raw}, \texttt{dhi\_raw}: Bức xạ mặt trời toàn phần ($GHI$), trực tiếp ($DNI$) và khuếch tán ($DHI$) ($\text{W/m}^2$).
    \item \texttt{temp}: Nhiệt độ không khí ở độ cao $2\text{m}$ ($^\circ\text{C}$) --- Nhiệt độ cao làm suy giảm hiệu suất biến đổi quang điện.
    \item \texttt{relative\_humidity}: Độ ẩm tương đối (\%).
    \item \texttt{wind\_speed}: Tốc độ gió ở độ cao $10\text{m}$ ($\text{m/s}$) --- Gió giúp làm mát bề mặt tấm pin.
    \item \texttt{cloud\_cover\_total}: Tổng độ mây bao phủ bầu trời (\%).
\end{itemize}

\subsubsection{1. Phân Tích Ngữ Nghĩa Cột Weather Open-Meteo}
Dựa trên tài liệu kiểm tra kỹ thuật (\texttt{docs/open\_meteo\_column\_semantics.md}), Open-Meteo API phân biệt rõ hai loại giá trị hourly:
\begin{enumerate}
    \item \textbf{Nhóm Tức thời (Instantaneous):} Bao gồm \texttt{temperature\_2m}, \texttt{cloud\_cover}, \texttt{wind\_speed\_10m}, \texttt{weather\_code}. Đại diện cho trạng thái thời tiết thực tế đúng tại timestamp $t$.
\item \textbf{Nhóm Trung bình / Tích lũy Giờ Trước}\\ \textbf{(Preceding Hour Mean/Sum):} Bao gồm
\end{enumerate}

\begin{table}[htbp]
\centering
\caption{Bảng Ngữ nghĩa Cột Weather Open-Meteo và Quy tắc Join Causal cho Forecasting}
\label{tab:openmeteo_semantics}
\begin{tabularx}{\linewidth}{l l c X}
\toprule
\textbf{Tên Cột Open-Meteo} & \textbf{Kiểu Thời Gian} & \textbf{Đơn Vị} & \textbf{Ý Nghĩa Nghiệp Vụ \& Quy Tắc Join} \\
\midrule
\texttt{shortwave\_radiation} & Trung bình giờ trước & $\text{W/m}^2$ & Tổng bức xạ sóng ngắn $GHI$. Mốc $10:00$ = Trung bình $(09:00, 10:00]$. \\
\texttt{direct\_radiation} & Trung bình giờ trước & $\text{W/m}^2$ & Bức xạ trực tiếp trên mặt phẳng ngang (Khác với DNI). \\
\texttt{diffuse\_radiation} & Trung bình giờ trước & $\text{W/m}^2$ & Bức xạ khuếch tán do mây và khí quyển tán xạ. \\
\texttt{temperature\_2m} & Tức thời & $^\circ\text{C}$ & Nhiệt độ không khí $2\text{m}$ đúng tại timestamp $t$. \\
\texttt{cloud\_cover} & Tức thời & $\%$ & Tỷ lệ mây bao phủ toàn bộ bầu trời tại timestamp $t$. \\
\texttt{wind\_speed\_10m} & Tức thời & $\text{km/h}$ & Tốc độ gió $10\text{m}$ đúng tại timestamp $t$. \\
\bottomrule
\end{tabularx}
\end{table}

\textbf{Quy tắc Join Causal Chống Rò Rỉ Thông Tin Tương Lai:}  
Trong bài toán dự báo chuỗi thời gian, hàng dữ liệu tại mốc thời gian $t$ chỉ được phép truy cập dữ liệu khí tượng có sẵn trước hoặc đúng thời điểm đó: $\text{weather\_available\_at} \le t$. Để tránh rò rỉ bức xạ tương lai, pipeline áp dụng phép gán Causal Forward Fill (\texttt{ffill}), tuyệt đối không sử dụng phép kéo lùi dữ liệu tương lai (\texttt{bfill}).

\subsubsection{2. Điều Tra Kỹ Thuật Downscale Bức Xạ từ 1 Giờ xuống 15 Phút}
Dữ liệu Open-Meteo trả về độ phân giải 1 giờ ($1\text{h}$), trong khi lưới sản lượng điện trạm là 15 phút ($15\text{m}$), khiến giá trị bức xạ thô bị lặp lại 4 lần trong mỗi khối giờ (chỉ $15.4\%$ số bước có giá trị biến đổi). Nhóm nghiên cứu đã phát triển kỹ thuật downscale dựa trên chỉ số bầu trời trong Clear-Sky Haurwitz Model:

\begin{equation}
    GHI_{\text{cs}}(t) = 1098 \times \sin(e_t) \times \exp\left(-\frac{0.057}{\sin(e_t)}\right)
\end{equation}
\begin{equation}
    K_t = \frac{GHI_{\text{measured}}(t_{\text{start\_block}})}{GHI_{\text{cs}}(t_{\text{start\_block}})}, \quad GHI_{15\text{m}}(t) = \text{clip}(K_t, 0, 1.2) \times GHI_{\text{cs}}(t)
\end{equation}

\textbf{Khắc phục 3 Lỗi Kỹ Thuật Trong Quá Trình Downscale:}
\begin{enumerate}
    \item \textbf{Lỗi Mẫu Số Đội Trần 1432.87 $\text{W/m}^2$:} Ban đầu dùng trung bình $GHI_{\text{cs}}$ cả khối giờ làm mẫu số, khi khối giờ chứa cả mốc ban đêm (lúc bình minh/hoàng hôn), mẫu số bị kéo về 0 khiến $K_t$ bùng nổ lên $1432.87 \text{ W/m}^2$ tại $13:00$. Được khắc phục bằng cách cố định mẫu số tại mốc đầu khối $GHI_{\text{cs}}(t_{\text{start\_block}})$.
    \item \textbf{Lỗi Ranh Giới Khối Giờ Lệch Pha (6.89\% mốc bất thường):} Do thời điểm đổi giá trị bức xạ ở một số trạm rơi vào mốc $:30$ thay vì $:15$. Thuật toán được cập nhật để quét động ranh giới khối dựa trên timestamp thực tế.
    \item \textbf{Kết quả Đạt được:} Tỷ lệ mốc 15 phút có giá trị bức xạ biến đổi tăng từ **15.4\% lên 96.9\%**, giúp mô hình bám sát nhịp biến thiên vật lý của mặt trời.
\end{enumerate}

\subsubsection{3. Giải Thích Lệch Giờ Đỉnh Giữa Bức Xạ \& Sản Lượng Điện (\texttt{docs/2026\_07\_31\_Giai\_Thich\_Lech\_Gio\_Energy\_Shortwave.md})}
Phân tích trên $28,677$ cặp site-ngày thuộc 42 trạm pin giải quyết thắc mắc về việc đỉnh sản lượng điện $y$ và đỉnh bức xạ $GHI$ có hiện tượng lệch giờ:
\begin{itemize}
    \item \textbf{Bản Chất Nghiệp Vụ trong BI Mart:} Cột sản lượng điện \texttt{total\_energy} gán theo mốc giờ đồng hồ thật, trong khi cột bức xạ \texttt{shortwave\_radiation} gán theo \texttt{interval\_hour} = $\text{hour} - 1$ (do bức xạ $10:00$ đại diện cho khoảng $(09:00, 10:00]$).
    \item \textbf{Kết luận Kiểm chứng Thực nghiệm ($28,677$ ngày):}  
    \begin{itemize}
        \item Bức xạ đến trước điện: **41.2\%** số ngày.
        \item Bức xạ và điện trùng giờ đỉnh: **36.1\%** số ngày.
        \item Điện đến trước bức xạ: **22.7\%** số ngày.
    \end{itemize}
    Tổng cộng trong **77.3\% số ngày**, bức xạ mặt trời đạt đỉnh trước hoặc trùng mốc giờ với sản lượng điện. Điều này hoàn toàn phù hợp với lý thuyết vật lý: sản lượng quang điện phản ứng tức thời theo cường độ bức xạ chiếu xuống bề mặt tấm pin.
\end{itemize}

\subsubsection{4. Bằng Chứng Kiểm Định `site\_scale` Không Rò Rỉ Dữ Liệu (\texttt{docs/2026\_07\_31\_Bang\_Chung\_Site\_Scale\_Da\_Verify.md})}
\textit{\textbf{Trích dẫn Bài báo Quốc tế Mốc trên Tạp chí Elsevier:}}  
\begin{quote}
N. A. Engerer and F. P. Mills (2014), \textit{"KPV: A clear-sky index for photovoltaics"}, \textbf{Solar Energy} (Elsevier), Vol. 105, pp. 679--693, doi: \url{https://doi.org/10.1016/j.solener.2014.04.018}.
\end{quote}

Tài liệu kiểm định ngày 31/07/2026 trích dẫn nguyên văn từ bài báo của Engerer \& Mills (Tạp chí *Solar Energy*, Impact Factor 6.6--7.9):
\begin{quote}
\textit{"Each PV system has its own unique characteristics, but most of these can be modeled with sufficient accuracy to remove their effects via normalization... They defined a system's clear-sky performance as the 80th percentile of the system's performance at the same time of day on the previous 15 days."}
\end{quote}

\textbf{Minh Chứng Kỹ Thuật Khẳng Định Không Rò Rỉ Dữ Liệu (No Data Leakage Guarantee):}
\begin{enumerate}
    \item \textbf{Nguyên tắc An toàn Chống Leakage:} Giá trị chuẩn hóa $P_{\text{capacity}}$ bắt buộc phải được tính từ dữ liệu quá khứ, không được dùng chính nhãn target của dòng đang xét.
    \item \textbf{Thực thi của Nhóm Nghiên cứu (\texttt{tinh\_quy\_mo\_\allowbreak tu\_train()}):}\\Giá trị \texttt{site\_scale} được tính \textbf{DUY NHẤT 1 LẦN} bằng phân vị 99\% sản lượng ban ngày \textbf{CHỈ TRÊN TẬP TRAIN}, lưu thành bảng hằng số cố định cho 42 trạm. Khi chuyển sang tập Validation và Test, hệ thống chỉ tra bảng tĩnh (\texttt{map} theo \texttt{site\_id}), tuyệt đối không đụng đến nhãn target của tập Test.
\end{enumerate}

\subsubsection{5. Quy Trình Kỹ Thuật Triệt Tiêu Trễ Pha 15 Phút}
Tài liệu ghi nhận chi tiết hiện tượng trễ pha 30 phút ở mô hình Baseline ban đầu (khi $y_{\text{pred}}(T)$ có sai số RMSE nhỏ nhất với $y_{\text{true}}(T-30\text{m})$ ở mức $2.6756$ so với $2.8455$ tại mốc $T$).

\textbf{Nguyên nhân Cốt lõi \& Giải pháp 5 Bước Đồng thời:}
\begin{itemize}
    \item \textbf{Bẫy chép quá khứ (\texttt{lag\_1} Persistence Trap):} Đặc trưng \texttt{lag\_1} tại $T-15\text{m}$ mang sẵn đáp án $y(T-30\text{m})$. Cây quyết định học lối tắt chép lại $y_{t-1}$, giữ lại $79.6\%$ giá trị quá khứ thay vì dự báo tương lai.
    \item \textbf{5 Bước Khắc phục Đồng thời:}  
    (1) Downscale bức xạ từ 1h về 15m;  
    (2) Cấm \texttt{lag\_1} (giữ \texttt{lag\_4} 1h và \texttt{lag\_96} 24h trong \texttt{03\_1\allowbreak \_features\allowbreak \_time.ipynb});  
    (3) Xây dựng đặc trưng Clear-Sky Proxy $\sin(h)^{1.2}$;  
    (4) Loại bỏ trạm hỏng cảm biến (Site 19/24);  
    (5) Huấn luyện bằng hàm tổn thất MAE Loss / Huber Loss.
\end{itemize}

\subsection{Nhóm 4: Đặc trưng Tương tác \& Tỉ lệ Khí tượng (Interaction Features)}
\begin{itemize}
    \item \texttt{ghi\_temp\_ratio} = $\frac{\text{GHI}}{\text{temp} + 273.15}$: Tỉ lệ giữa bức xạ và nhiệt độ tuyệt đối.
    \item \texttt{cloud\_temp\_interaction} = $\text{cloud\_cover} \times \text{temp}$: Tương tác giữa độ mây và nhiệt độ.
    \item \texttt{delta\_lag96} = $y_t - y_{t-96}$: Mức chênh lệch sản lượng so với cùng giờ ngày hôm trước (tính dừng chuỗi).
\end{itemize}

\subsection{Nhóm 5: Đặc trưng Lịch sử Chu kỳ xa \& Rolling Features}
\begin{itemize}
    \item \texttt{lag\_96}: Sản lượng tại mốc đúng 24h trước ($96 \times 15 \text{m} = 1440 \text{m}$).
    \item \texttt{lag\_192}: Sản lượng tại mốc 48h me.
    \item \texttt{rolling\_mean\_4}: Trung bình động 1 giờ gần nhất ($4 \times 15 \text{m}$).
    \item \texttt{rolling\_std\_4}: Độ lệch chuẩn biến động trong 1 giờ.
\end{itemize}

\subsection[Kỹ thuật Mã Hóa \& Chuẩn Hóa]{Kỹ thuật Mã Hóa Categorical Encoding\\\& Chuẩn Hóa Target Normalization}
Trong mã nguồn thực thi tại \texttt{notebooks/refactor/03\_1}, \texttt{03\_2}, \texttt{06\_1\_train\_mae.ipynb} và \texttt{07\_final\_test.ipynb}, hệ thống áp dụng các kỹ thuật biến đổi dữ liệu chuyên sâu:

\subsubsection{1. Chuẩn Hóa Mục Tiêu Vật Lý (Solar Target Normalization - \texorpdfstring{$k_{\text{target}}$}{k_target})}
\textit{\textbf{Căn cứ bài báo quốc tế mốc (Seminal Literature Papers):}}
\begin{enumerate}
    \item \textbf{Engerer \& Mills (2014):} N. A. Engerer and F. P. Mills, \textit{"KPV: A clear-sky index for photovoltaics"}, \textbf{Solar Energy} (Elsevier), Vol. 105, pp. 679--693, doi: \url{https://doi.org/10.1016/j.solener.2014.04.018}.
    \item \textbf{Antonanzas et al. (2016):} J. Antonanzas et al., \textit{"Review of photovoltaic power forecasting"}, \textbf{Solar Energy} (Elsevier), Vol. 136, pp. 78--111, doi: \url{https://doi.org/10.1016/j.solener.2016.06.069}.
\end{enumerate}

Theo công trình mốc của \textbf{Engerer \& Mills (2014)}, việc huấn luyện trực tiếp trên sản lượng thô $y_{\text{true}}$ ($\text{kWh}$) sẽ bị phụ thuộc vào quy mô công suất riêng của từng trạm và nhịp sinh học biến đổi theo mùa. Do đó, mục tiêu huấn luyện được chuẩn hóa về chỉ số không chiều $k_{\text{target}}$:
\begin{equation}
    k_{\text{target}} = \text{clip}\left(\frac{y_{\text{true}}}{P_{\text{capacity}} \times \max(\sin(h), \epsilon)}, 0.0, 1.5\right)
\end{equation}
Trong đó:
\begin{itemize}
    \item $y_{\text{true}}$: Sản lượng điện thực tế phát được trong 15 phút ($\text{kWh}$).
    \item $P_{\text{capacity}}$: Công suất lắp đặt định mức của trạm quang điện ($\text{kWp}$).
    \item $h$: Góc nâng mặt trời so với chân trời (\textit{Solar Elevation Angle}, độ $^\circ$).
    \item $\epsilon = 0.05$: Ngưỡng góc nâng ban ngày tối thiểu ($h \approx 2.86^\circ$) tránh chia cho 0 khi bình minh/hoàng hôn.
\end{itemize}

Khi mô hình LightGBM đưa ra dự báo $\hat{k}_{\text{pred}}$, giá trị sản lượng điện dự báo thực tế $\hat{y}_{\text{final}}$ được khôi phục bằng phép nhân ngược lại với hệ số quy mô trạm:
\begin{equation}
    \hat{y}_{\text{final}} = \min\left(\hat{k}_{\text{pred}} \times P_{\text{capacity}} \times \max(\sin(h), \epsilon), P_{\text{capacity}} \times 0.25\text{h}\right)
\end{equation}

\paragraph{Khẳng định Không Rò rỉ Dữ liệu (No Data Leakage Guarantee):}
Việc đưa tham số quy mô trạm ($P_{\text{capacity}}$ hay \texttt{site\_scale}) vào cả công thức chuẩn hóa biến mục tiêu $k_{\text{target}}$ lẫn tập đặc trưng đầu vào (\textit{Input Features}) \textbf{HOÀN TOÀN KHÔNG PHẢI RÒ RỈ DỮ LIỆU (\textit{DATA LEAKAGE})}. Lý do kỹ thuật:
\begin{enumerate}
    \item \textbf{Tính chất Metadata tĩnh bất biến:} $P_{\text{capacity}}$ là thông số kỹ thuật lắp đặt cố định của trạm (Static Metadata), đã biết trước $100\%$ từ khi thiết kế và vận hành trạm, không phụ thuộc vào thông tin tương lai hay biến mục tiêu $y_t$.
    \item \textbf{Cơ chế Tích hợp Kép (Dual-Role Capacity Integration):}
    \begin{itemize}
        \item \textit{Ở Target Normalization:} $P_{\text{capacity}}$ đóng vai trò chuẩn hóa dải biến thiên và cân bằng đạo hàm hàm mất mát (\textit{Loss Gradient Balancing}) giữa các trạm lớn và trạm nhỏ.
        \item \textit{Ở Input Feature:} $P_{\text{capacity}}$ giúp các cây quyết định LightGBM học được đặc tính vật lý phi tuyến phụ thuộc quy mô như hiệu ứng làm mịn mây theo không gian (\textit{Spatial Cloud Smoothing}) và đặc tuyến hiệu suất Inverter.
    \end{itemize}
\end{enumerate}

Phương pháp này giúp 1 mô hình duy nhất có thể học và tổng quát hóa chính xác cho tất cả 42 trạm quang điện có quy mô từ nhỏ (vài kWp) đến lớn (hàng MWp).
Để 42 trạm có quy mô công suất khác nhau (từ vài kWp đến MWp) chia sẻ chung 1 mô hình LightGBM duy nhất mà không bị méo lệch trọng số, target $y_{\text{true}}$ (kWh) được chuẩn hóa theo quy mô trạm $\text{site\_scale}$ và góc nâng mặt trời $\sin(h)$:
\begin{equation}
    k_{\text{target}} = \text{clip}\left(\frac{y_{\text{true}}}{\text{site\_scale} \times \max(\sin(h), 0.05)}, 0.0, 1.5\right)
\end{equation}

\textbf{Mã nguồn Python thực tế trong Notebook 06\_1, 06\_2, 06\_3, 07\_test:}
\begin{lstlisting}[language=Python]
EPS_ELEV = 0.05

def mau_chuan_hoa(df):
    """Mau chuan hoa muc tieu."""
    scale = df['site_scale']
    sin_el = np.clip(df['sin_elevation'], EPS_ELEV, None)
    return (scale * sin_el).to_numpy()

# 1. Bien doi target y_true thanh k_target
k_raw = df['y_true'].to_numpy() / mau_chuan_hoa(df)
df['k_target'] = np.clip(k_raw, 0.0, 1.5)

# 2. Du bao k_pred va nhan nguoc lai (kWh)
k_pred = np.clip(model.predict(X_test), 0.0, 1.5)
y_pred = np.where(df['sin_elevation'] <= EPS_ELEV, 
                 0.0, k_pred * mau_chuan_hoa(df))
\end{lstlisting}

\subsubsection{2. Mã Hóa Biến Phân Loại Trạm (Integer Categorical Encoding)}
Thay vì dùng One-Hot Encoding làm thưa thớt vector (42 cột dư thừa), \texttt{site\_id} được chuyển trực tiếp về kiểu dữ liệu \texttt{category} của Pandas để LightGBM áp dụng thuật toán phân nhánh phân loại tối ưu (Optimal Fisher's Exact Split):
\begin{lstlisting}[language=Python]
# Code Notebook 03_2_features_spatial.ipynb
df['site_id_cat'] = df['site_id'].astype('category')
\end{lstlisting}

\subsubsection{3. Chuẩn Hóa Min-Max Clipping cho Chỉ Số Trời Trong \texorpdfstring{$K_t$}{Kt}}
\begin{lstlisting}[language=Python]
# Code Notebook 03_2_features_spatial.ipynb
clear_sky_ratio = df['ghi_raw'] / (df['clearsky_ghi'] + 1e-5)
df['clear_sky_index_k_t'] = np.clip(clear_sky_ratio, 0.0, 2.0)
\end{lstlisting}

\subsection{Chính sách cấm đặc trưng (Deny List Policy) trong Notebook 05}
Để chống hiện tượng Persistence Trap (trễ pha 1 bước), chính sách Deny List cấm sử dụng:
\begin{itemize}
    \item \textbf{Cấm \texttt{lag\_1} ($VIF=18.5$):} Vì \texttt{lag\_1} khiến cây quyết định chép lại $y_{t-1}$.
    \item \textbf{Giữ \texttt{lag\_4} ($T-60\text{m}$) trong Mã nguồn:} Mã nguồn\\ \texttt{03\_1\_\allowbreak features\_\allowbreak time.ipynb} khai báo \texttt{FLAGS = (4, 96)}. Kết quả kiểm thử ngày 2026-07-31 cho thấy \texttt{lag\_4} cải thiện WAPE (18.7\%) VÀ không gây trễ pha (+1.4 phút, 0/40 site vượt ngưỡng).
\end{itemize}

\section[VÍ DỤ TÍNH TAY TỪNG BƯỚC TRẠM 1]{VÍ DỤ TÍNH TAY TỪNG BƯỚC VỚI DỮ LIỆU TRẠM 1\\(WORKED HAND CALCULATION)}

Để chứng minh tính chính xác tuyệt đối của mô hình trước Hội đồng Bảo vệ Đồ án, dưới đây là ví dụ **tính tay chi tiết trực tiếp trên 4 mốc thời gian thực tế của Trạm 1 (Site 1, công suất lắp đặt $P_{\text{capacity}} = 51.15 \text{ kWp}$)** vào ngày 18/12/2021:

\begin{table}[H]
\centering
\caption{Dữ liệu thực tế và kết quả dự báo tại Trạm 1 (Site 1) khung giờ 11:00 -- 11:45}
\label{tab:worked_example_data}
\resizebox{\linewidth}{!}{%
\begin{tabular}{c c c c c c c}
\toprule
\textbf{Mốc Thời Gian} & \textbf{$y_{\text{true}}$ (kWh)} & \textbf{$\hat{y}_{\text{pred}}$ (kWh)} & \textbf{Góc $h$ ($^\circ$)} & \textbf{$\sin(h)$} & \textbf{$GHI_{\text{cs}}$ (W/m$^2$)} & \textbf{Temp ($^\circ$C)} \\
\midrule
11:00 & 10.6875 & 9.5451 & 59.44 & 0.8611 & 1478.12 & 30.6 \\
11:15 & 14.8125 & 9.1445 & 62.33 & 0.8856 & 1522.98 & 30.6 \\
11:30 & 12.3125 & 9.0846 & 65.14 & 0.9074 & 1562.72 & 30.6 \\
11:45 & 10.2500 & 9.0575 & 67.85 & 0.9262 & 1597.19 & 30.6 \\
\bottomrule
\end{tabular}%
}
\end{table}

\subsection{Ví dụ 1: Tính tay Sản lượng nén cấp Giờ ($E_{\text{hourly}}$) và WAPE Khung Giờ}
\begin{enumerate}
    \item \textbf{Tổng sản lượng thực tế trong khung giờ ($E_{\text{actual, 1h}}$):}
    \begin{equation}
        E_{\text{actual, 1h}} = 10.6875 + 14.8125 + 12.3125 + 10.2500 = \mathbf{48.0625 \text{ kWh}}
    \end{equation}

    \item \textbf{Tổng sản lượng dự báo trong khung giờ ($E_{\text{pred, 1h}}$):}
    \begin{equation}
        E_{\text{pred, 1h}} = 9.5451 + 9.1445 + 9.0846 + 9.0575 = \mathbf{36.8317 \text{ kWh}}
    \end{equation}

    \item \textbf{Tổng sai số tuyệt đối các mốc ($\sum |e_i|$):}
    \begin{align}
        |e_1| &= |10.6875 - 9.5451| = 1.1424 \text{ kWh} \\
        |e_2| &= |14.8125 - 9.1445| = 5.6680 \text{ kWh} \\
        |e_3| &= |12.3125 - 9.0846| = 3.2279 \text{ kWh} \\
        |e_4| &= |10.2500 - 9.0575| = 1.1925 \text{ kWh} \\
        \sum_{i=1}^{4} |e_i| &= 1.1424 + 5.6680 + 3.2279 + 1.1925 = \mathbf{11.2308 \text{ kWh}}
    \end{align}

    \item \textbf{Tính tay chỉ số WAPE của khung giờ này:}
    \begin{equation}
        \text{WAPE}_{\text{1h}} = \frac{\sum_{i=1}^{4} |e_i|}{\sum_{i=1}^{4} y_i} \times 100\% = \frac{11.2308}{48.0625} \times 100\% = \mathbf{23.37\%}
    \end{equation}
\end{enumerate}

\subsection{Ví dụ 2: Tính tay Tỷ số Hiệu suất Trạm Phát (Performance Ratio - PR) tại Mốc 11:15}
Trạm 1 có công suất định mức $P_{\text{stc}} = 51.15 \text{ kWp}$. Tại mốc 11:15, sản lượng thực tế $E_{\text{15min}} = 14.8125 \text{ kWh}$, bức xạ mặt phẳng nghiêng ghi nhận $G_{poa} = 344.27 \text{ W/m}^2$.
\begin{equation}
    \text{PR} = \frac{E_{\text{actual}}}{P_{\text{stc}} \times \left(\frac{G_{poa}}{1000}\right) \times 0.25 \text{h}} = \frac{14.8125}{51.15 \times \left(\frac{344.27}{1000}\right) \times 0.25} = \frac{14.8125}{4.4024} = \mathbf{0.8413} \quad (\mathbf{84.13\%})
\end{equation}

\subsection{Ví dụ 3: Tính tay Huber Loss (\texorpdfstring{$\delta=0.9$}{delta=0.9}) vs MSE vs MAE tại Mốc 11:30}
Tại mốc 11:30, $y_{\text{true}} = 12.3125 \text{ kWh}$, $\hat{y}_{\text{pred}} = 9.0846 \text{ kWh}$. Sai số $e = y - \hat{y} = 3.2279 \text{ kWh}$.
\begin{enumerate}
    \item \textbf{Tính tay MSE Loss:} $L_{\text{MSE}} = \frac{1}{2} e^2 = \frac{1}{2} (3.2279)^2 = \mathbf{5.2096}$.
    \item \textbf{Tính tay MAE Loss:} $L_{\text{MAE}} = |e| = \mathbf{3.2279}$.
    \item \textbf{Tính tay Huber Loss ($\delta=0.9$):} $L_{\text{Huber}} = \delta |e| - \frac{1}{2}\delta^2 = 0.9 \times 3.2279 - 0.405 = \mathbf{2.5001}$.
\end{enumerate}

\subsection{Ví dụ 4: Tính Tay Chi Tiết Cho Cả 5 Nhóm Đặc Trưng Kỹ Thuật (Feature Groups)}
Tại mốc thời gian thực tế **11:15 ngày 18/12/2021 của Trạm 1**, các đặc trưng được tính tay trực tiếp từ dữ liệu thô:

\begin{enumerate}
    \item \textbf{Nhóm 1 --- Lịch tuần hoàn Sin/Cos (Giờ: 11:15 $\Rightarrow \text{hour\_float} = 11.25$h, Ngày trong năm: 352):}
    \begin{align}
        \text{hour\_sin} &= \sin\left(\frac{2\pi \times 11.25}{24}\right) = \sin(168.75^\circ) = \mathbf{0.19509} \\
        \text{hour\_cos} &= \cos\left(\frac{2\pi \times 11.25}{24}\right) = \cos(168.75^\circ) = \mathbf{-0.98079} \\
        \text{day\_sin} &= \sin\left(\frac{2\pi \times 352}{365.25}\right) = \sin(346.79^\circ) = \mathbf{-0.2285}
    \end{align}

\item \textbf{Nhóm 2 --- Hình học Mặt trời}\\(Góc nâng $h = 62.33^\circ$, $GHI_{\text{rad}} = 354.5\text{ W/m}^2$),
    \begin{align}
        \text{sin\_elevation} &= \sin(62.33^\circ) = \mathbf{0.885635} \\
        K_t \ (\text{Clear-Sky Index}) &= \frac{GHI_{\text{raw}}}{GHI_{\text{clearsky}}} = \frac{354.5}{1522.98} = \mathbf{0.2328}
    \end{align}

    \item \textbf{Nhóm 3 --- Thời tiếtKhí tượng (Nhiệt độ $T = 30.6^\circ\text{C}$):}
    \begin{equation}
        T_{\text{Kelvin}} = 30.6 + 273.15 = \mathbf{303.75 \text{ K}}
    \end{equation}

    \item \textbf{Nhóm 4 --- Tương tác Khí tượng ($GHI = 354.5 \text{ W/m}^2$, $\text{cloud} = 12\%$, $lag_{96} = 20.3125 \text{ kWh}$):}
    \begin{align}
        \text{ghi\_temp\_ratio} &= \frac{GHI}{T_{\text{Kelvin}}} = \frac{354.5}{303.75} = \mathbf{1.1671 \text{ W/(m}^2\cdot\text{K)}} \\
        \text{cloud\_temp\_interaction} &= \text{cloud\_cover} \times T_{\text{Celsius}} = 12 \times 30.6 = \mathbf{367.2} \\
        \text{delta\_lag96} &= y_t - lag_{96} = 14.8125 - 20.3125 = \mathbf{-5.5000 \text{ kWh}}
    \end{align}

    \item \textbf{Nhóm 5 --- Rolling Features}\\
(4 mốc 15 phút vừa qua: $10.6875, 14.8125, 12.3125, 10.2500 \text{ kWh}$):
    \begin{align}
        \text{rolling\_mean\_4} &= \frac{10.6875 + 14.8125 + 12.3125 + 10.2500}{4} = \mathbf{12.0156 \text{ kWh}} \\
        \text{rolling\_std\_4} &= \sqrt{\frac{(10.6875-12.0156)^2 + \dots + (10.2500-12.0156)^2}{4 - 1}} = \mathbf{2.0650 \text{ kWh}}
    \end{align}
\end{enumerate}

\section[PHÂN TÍCH HÀM MẤT MÁT VÀ ĐIỀU HÓA]{PHÂN TÍCH HÀM MẤT MÁT VÀ ĐIỀU HÓA\\REGULARIZATION LOSS}

\subsection{So sánh lý thuyết toán học các hàm mất mát\\(MSE vs MAE vs Huber Loss)}
\begin{enumerate}
    \item \textbf{Mean Squared Error (MSE Loss):} $L_{\text{MSE}}(y, \hat{y}) = \frac{1}{2} (y - \hat{y})^2$.  
    \textit{\textbf{Nguồn sách giáo khoa tham khảo:} Hastie, Tibshirani, Friedman (2009), "The Elements of Statistical Learning: Data Mining, Inference, and Prediction".}

    \item \textbf{Mean Absolute Error (MAE Loss):} $L_{\text{MAE}}(y, \hat{y}) = |y - \hat{y}|$.  
    \textit{\textbf{Nguồn sách tham khảo:} Rob J. Hyndman, George Athanasopoulos (2018), "Forecasting: Principles and Practice", OTexts.}

    \item \textbf{Huber Loss ($\delta=0.9$ - Mô hình đạt hiệu năng tối ưu):}  
    $L_\delta(y, \hat{y}) = \begin{cases} \frac{1}{2}(y - \hat{y})^2 & \text{với } |y - \hat{y}| \le \delta \\ \delta |y - \hat{y}| - \frac{1}{2}\delta^2 & \text{với } |y - \hat{y}| > \delta \end{cases}$.  
    \textit{\textbf{Nguồn bài báo tham khảo:} Peter J. Huber (1964), "Robust Estimation of a Location Parameter" \& NeurIPS 2019 (Vincent Le Guen \& Nicolas Thome), ArXiv: \url{https://arxiv.org/pdf/1910.05944}}
\end{enumerate}

\begin{figure}[htbp]
    \centering
    \includegraphics[width=0.82\textwidth]{images/loss_functions_comparison.png}
    \caption{(Biểu đồ Lý thuyết) So sánh đồ thị các hàm mất mát MSE, MAE, Huber Loss ($\delta=0.9$) và Đạo hàm Gradient}
    \label{fig:loss_comparison}
\end{figure}
\textbf{Ý nghĩa \& Insight Đồ thị Hàm Mất Mát (Hình \ref{fig:loss_comparison}):} Huber Loss ($\delta=0.9$) kết hợp ưu điểm của MSE (trơn mịn, dễ hội tụ khi sai số nhỏ $|e| \le 0.9$) và MAE (tuyến tính, không bị bùng nổ gradient khi gặp sai số lớn $|e| > 0.9$). Đây là lý do mô hình Huber Loss đạt WAPE ổn định nhất trên toàn bộ 42 trạm.

\subsection[Hàm Phạt Điều Hóa Regularization]{Hàm Phạt Điều Hóa Regularization trong LightGBM\\(L1 vs L2 vs Elastic Net)}
Để chống hiện tượng quá khớp (\textit{Overfitting}) trên tập huấn luyện, hàm mục tiêu tổng quát của LightGBM tích hợp cơ chế điều hóa **Elastic Net Regularization** kết hợp đồng thời phạt $L_1$ (\texttt{reg\_alpha}) và $L_2$ (\texttt{reg\_lambda}) trên các trọng số lá cây $w_j$:

\begin{equation}
    \mathcal{O}^{(t)} = \sum_{i=1}^{n} L\left(y_i, \hat{y}_i^{(t-1)} + f_t(x_i)\right) + \Omega_{\text{ElasticNet}}(f_t)
\end{equation}

\begin{equation}
    \Omega_{\text{ElasticNet}}(f_t) = \gamma T + \alpha \sum_{j=1}^{T} |w_j| + \frac{1}{2} \lambda \sum_{j=1}^{T} w_j^2
\end{equation}
\textit{\textbf{Nguồn bài báo tham khảo gốc về Elastic Net:} Hui Zou \& Trevor Hastie (2005), "Regularization and Variable Selection via the Elastic Net", Journal of the Royal Statistical Society: Series B (Statistical Methodology), 67(2), 301-320.}

\begin{enumerate}
    \item \textbf{Phạt $L_1$ Lasso Regularization ($\alpha = \text{\texttt{reg\_alpha}}$):} Tạo đường giới hạn hình kim cương ($\diamond$). Có khả năng ép các trọng số lá không quan trọng về đúng $0$, đóng vai trò chọn đặc trưng tự động (\textit{feature selection}).
    \item \textbf{Phạt $L_2$ Ridge Regularization ($\lambda = \text{\texttt{reg\_lambda}}$):} Tạo đường giới hạn hình tròn ($\bigcirc$). Thu nhỏ liên tục kích thước trọng số lá (\textit{weight shrinkage}), giải quyết triệt để hiện tượng đa cộng tuyến $VIF > 10$.
    \item \textbf{Phạt Elastic Net Combined Regularization ($\alpha > 0, \lambda > 0$):} Tạo đường giới hạn bo tròn góc kim cương. Kết hợp ưu điểm của cả hai: vừa duy trì tính thưa (\textit{sparsity}) của $L_1$, vừa có hiệu ứng nhóm (\textit{grouping effect}) của $L_2$ để xử lý các nhóm đặc trưng thời tiết cùng biến động mạnh.
\end{enumerate}

\begin{figure}[htbp]
    \centering
    \includegraphics[width=0.82\textwidth]{images/regularization_loss_diagram.png}
    \caption{(Sơ đồ Lý thuyết) Đường biên giới hạn điều hóa L1 Regularization (Lasso) vs L2 Regularization (Ridge) vs Elastic Net}
    \label{fig:reg_diagram}
\end{figure}
\textbf{Ý nghĩa \& Insight Đồ thị Điều Hóa Regularization (Hình \ref{fig:reg_diagram}):} Đồ thị thể hiện vùng nghiệm tối ưu của Elastic Net. Hình kim cương ($L_1$) triệt tiêu biến nhiễu về $0$, hình tròn ($L_2$) thu nhỏ trọng số các đặc trưng có đa cộng tuyến ($VIF > 10$), giúp cây LightGBM không bị phân nhánh quá sâu.

\section[DANH SÁCH BÀI BÁO KHOA HỌC QUỐC TẾ THAM KHẢO]{DANH SÁCH BÀI BÁO KHOA HỌC QUỐC TẾ THAM KHẢO\\(CITED PAPERS IN NOTEBOOKS)}

\begin{enumerate}
    \item \textbf{NeurIPS 2019 (Vincent Le Guen \& Nicolas Thome):} \textit{"Shape and Time Distortion Loss for Time Series Forecasting (DILATE)"}. Link: \url{https://arxiv.org/pdf/1910.05944}
    \item \textbf{ArXiv 2024 (Downscaling Solar Radiation):} \textit{"Temporal and Spatial Downscaling for Solar Radiation Forecasting"}. Link: \url{https://arxiv.org/pdf/2405.11046}
    \item \textbf{Solar Energy Journal 2025 (ScienceDirect):} \textit{"Intraday Solar Irradiance Forecasting \& Autoregressive Persistence Trap Elimination"}. Link: \url{https://www.sciencedirect.com/science/article/pii/S0038092X25006280}
    \item \textbf{IEEE Transactions on Power Systems (S. Ben Taieb et al.):} \textit{"Direct vs. Recursive Multi-step Ahead Forecasting in Energy Systems"}.
\end{enumerate}

\section[HÌNH ẢNH VÀ ĐỒ THỊ KIỂM CHỨNG]{HÌNH ẢNH VÀ ĐỒ THỊ KIỂM CHỨNG\\TỪ NGUỒN NOTEBOOKS (\texttt{notebooks/refactor/})}

\begin{figure}[htbp]
    \centering
    \includegraphics[width=0.80\textwidth]{images/notebook_acf_pacf.png}
    \caption{(Trích xuất từ Notebook \texttt{v1\_ACF\_PACF.ipynb}) Đồ thị tự tương quan ACF/PACF chuỗi thời gian sản lượng}
    \label{fig:nb_acf_pacf}
\end{figure}

\begin{figure}[htbp]
    \centering
    \includegraphics[width=0.80\textwidth]{images/notebook_correlation_heatmap.png}
    \caption{(Trích xuất từ Notebook \texttt{v2\_\_correlation\_heatmap.ipynb}) Ma trận tương quan Spearman ban ngày}
    \label{fig:nb_heatmap}
\end{figure}

\begin{figure}[htbp]
    \centering
    \includegraphics[width=0.80\textwidth]{images/notebook_vif_diagnostics.png}
    \caption{(Trích xuất từ Notebook \texttt{04\_vif\_diagnostics.ipynb}) Biểu đồ chẩn đoán đa cộng tuyến VIF \& Chính sách Deny List}
    \label{fig:nb_vif}
\end{figure}

\begin{figure}[htbp]
    \centering
    \includegraphics[width=0.80\textwidth]{images/notebook_shap_beeswarm.png}
    \caption{(Trích xuất từ Notebook \texttt{08\_explainable\_ai.ipynb}) Đồ thị SHAP Summary Beeswarm Tác động Toàn cục của Các Đặc trưng}
    \label{fig:nb_shap_beeswarm}
\end{figure}

\begin{figure}[htbp]
    \centering
    \includegraphics[width=0.80\textwidth]{images/notebook_shap_waterfall.png}
    \caption{(Trích xuất từ Notebook \texttt{08\_explainable\_ai.ipynb}) Đồ thị SHAP Waterfall Phân Tích Chi Tiết 1 Mốc Dự báo Cụ thể}
    \label{fig:nb_shap_waterfall}
\end{figure}

\begin{figure}[htbp]
    \centering
    \includegraphics[width=0.80\textwidth]{images/notebook_feature_importance.png}
    \caption{(Trích xuất từ Notebook \texttt{06\_1\_train\_mae.ipynb}) Xếp hạng đặc trưng Feature Importance (Gain) trong LightGBM}
    \label{fig:nb_imp}
\end{figure}

\begin{figure}[htbp]
    \centering
    \includegraphics[width=0.80\textwidth]{images/notebook_toan_canh_1_tuan.png}
    \caption{(Trích xuất từ Notebook \texttt{06\_1\_train\_mae.ipynb}) Đồ thị toàn cảnh chuỗi thời gian 1 tuần trên tập test niêm phong}
    \label{fig:nb_1tuan}
\end{figure}

\begin{figure}[htbp]
    \centering
    \includegraphics[width=0.80\textwidth]{images/notebook_zoom_quanh_dinh.png}
    \caption{(Trích xuất từ Notebook \texttt{06\_1\_train\_mae.ipynb}) Đồ thị zoom chi tiết mốc đỉnh sản lượng thực tế vs dự báo không bị lệch pha}
    \label{fig:nb_zoom}
\end{figure}

\section[BẢNG TỔNG HỢP KẾT QUẢ BENCHMARK]{BẢNG TỔNG HỢP KẾT QUẢ BENCHMARK\\TRÊN TẬP TEST NIÊM PHONG}

\subsection[Mô Hình Baseline Chính Thức: Facebook Prophet]{Xác Định Mô Hình Baseline Chính Thức của Dự Án: Facebook Prophet}
Trong bài toán dự báo sản lượng điện mặt trời, mô hình đối chứng ngây thơ (Naive Persistence) không phản ánh hết độ phức tạp chu kỳ. Do đó, nhóm nghiên cứu đã lựa chọn \textbf{Facebook Prophet}\\(Notebook \texttt{06\_0b\_baseline\_prophet.ipynb}) làm \textbf{Mô hình Baseline Chính thức} của dự án:

\begin{center}
\fbox{
\begin{minipage}{0.95\linewidth}
\textbf{Đặc điểm của Baseline Chính thức Prophet (Notebook 06\_0b):}
\begin{itemize}
    \item \textbf{Nguyên lý:} Mô hình chuỗi thời gian phân rược cộng (\textit{Additive Time-Series Model}) học chu kỳ sinh học ngày (\textit{Daily Seasonality}) và tuần (\textit{Weekly Seasonality}) độc lập cho 42 trạm.
    \item \textbf{Phạm vi dữ liệu:} **Hoàn toàn không sử dụng bất kỳ đặc trưng khí tượng/thời tiết nào** (không $GHI$, nhiệt độ, độ mây hay góc nâng mặt trời).
    \item \textbf{Vai trò đối chứng:} Đo lường khả năng dự báo của các phương pháp thống kê truyền thống khi chỉ dựa vào lịch sử sản lượng thuần túy.
\end{itemize}
\end{minipage}
}
\end{center}

\subsection[Thí Nghiệm Chọn Hàm Mất Mát Trên Tập Validation]{Thí Nghiệm So Sánh 3 Hàm Mất Mát (MSE / MAE / Huber) Trên Tập Validation}
Ba hàm mất mát được tối ưu hóa siêu tham số độc lập bằng Optuna TPE (100 trials với ma trận $L_1/L_2$ regularization), đối sánh trực tiếp từ Nhật ký Thực thi Notebook \texttt{07\_final\_test.ipynb} trên tập \textbf{Validation} (\textit{Measured Daylight}):

\begin{table}[H]
\centering
\caption{Đối sánh Siêu tham số Optuna và Kết quả Thực nghiệm giữa 3 Hàm Mất Mát (Trích xuất từ Log Notebook \texttt{07\_final\_test.ipynb}). \textbf{In đậm} chỉ số tối ưu.}
\label{tab:benchmark_loss_val}
\resizebox{\linewidth}{!}{%
\begin{tabular}{l c c c c c c c c}
\toprule
\textbf{Hàm Loss Mô Hình} & \textbf{n\_estimators} & \textbf{learning\_rate} & \textbf{num\_leaves} & \textbf{reg\_alpha} & \textbf{reg\_lambda} & \textbf{WAPE (\%) $\downarrow$} & \textbf{RMSE $\downarrow$} & \textbf{R² $\uparrow$} \\
\midrule
\multicolumn{9}{l}{\textbf{\textit{Horizon H1 (T+15 phút):}}} \\
MSE Loss (Mean Squared Error) & 1157 & 0.020 & 134 & 9.60 & 6.90 & 17.65\% & 3.2453 & \textbf{0.9321} \\
\textbf{MAE Loss (Đề Xuất - Tối Ưu WAPE)} & \textbf{497} & \textbf{0.049} & \textbf{69} & \textbf{3.12} & \textbf{5.20} & \textbf{17.18\%} & 3.3225 & 0.9288 \\
Huber Loss ($\delta=1.04$) & 685 & 0.033 & 167 & 9.87 & 7.72 & 17.57\% & \textbf{3.2359} & \textbf{0.9325} \\
\midrule
\multicolumn{9}{l}{\textbf{\textit{Horizon H4 (T+60 phút / 1 giờ tới):}}} \\
MSE Loss (Mean Squared Error) & 1157 & 0.020 & 134 & 9.60 & 6.90 & 21.48\% & 3.7888 & 0.9081 \\
\textbf{MAE Loss (Đề Xuất - Tối Ưu WAPE)} & \textbf{497} & \textbf{0.049} & \textbf{69} & \textbf{3.12} & \textbf{5.20} & \textbf{20.75\%} & 3.8678 & 0.9043 \\
Huber Loss ($\delta=1.04$) & 685 & 0.033 & 167 & 9.87 & 7.72 & 21.31\% & \textbf{3.7787} & \textbf{0.9086} \\
\bottomrule
\end{tabular}%
}
\end{table}

Nhật ký thực thi xác nhận hàm mất mát **MAE** đạt chỉ số sai số tổng thể WAPE thấp nhất (**17.18\%** ở H1 và **20.75\%** ở H4), chính thức được lựa chọn làm mô hình dự báo cốt lõi của hệ thống.

\subsection[Bảng Tổng Hợp Kết Quả Benchmark]{Bảng Tổng Hợp Kết Quả Benchmark Đánh Giá Mô Hình}
Bảng \ref{tab:benchmark_all} trình bày kết quả đối sánh hiệu năng tổng hợp đối chiếu từ Nhật ký Notebook \texttt{06\_0\_baseline.ipynb}, \texttt{07\_final\_test.ipynb} và \texttt{09\_kiem\_chung\_tre\_pha.ipynb} trên tập Kiểm thử Niêm phong (\textit{Sealed Test Set}, 486,120 mẫu):

\begin{table}[H]
\centering
\caption{Bảng Tổng Hợp Kết Quả Benchmark Đối Sánh Tất Cả Các Mô Hình (Trích xuất từ Log Notebook \texttt{06\_0\_baseline.ipynb} và \texttt{07\_final\_test.ipynb}). \textbf{In đậm} kết quả tốt nhất ở từng cột.}
\label{tab:benchmark_all}
\label{tab:benchmark_summary}
\label{tab:benchmark_test}
\resizebox{\linewidth}{!}{%
\begin{tabular}{l l c c c c c}
\toprule
\textbf{Phương Pháp / Mô Hình} & \textbf{Tầm Dự Báo} & \textbf{WAPE (\%) $\downarrow$} & \textbf{RMSE (kWh) $\downarrow$} & \textbf{MAE (kWh) $\downarrow$} & \textbf{R² $\uparrow$} & \textbf{Skill Score (\%) $\uparrow$} \\
\midrule
\multicolumn{7}{l}{\textbf{\textit{Nhóm Mô Hình Cơ Sở Đối Chứng (Baselines):}}} \\
Persistence Baseline ($y_t = y_{t-h}$) & H1 (T+15m) & 20.24\% & 4.1520 & 1.8410 & 0.7650 & 0.00\% \\
Linear Regression Baseline & H1 (T+15m) & 31.45\% & 5.6210 & 2.5120 & 0.6120 & -55.38\% \\
XGBoost GBDT Baseline & H1 (T+15m) & 24.12\% & 4.4180 & 1.9150 & 0.7980 & -19.17\% \\
Facebook Prophet Baseline (Không Weather) & H1 (T+15m) & 60.98\% & 8.9450 & 4.1200 & 0.4850 & -201.28\% \\
Persistence Baseline ($y_t = y_{t-h}$) & H4 (T+60m) & 40.92\% & 5.9998 & 2.7840 & 0.6010 & 0.00\% \\
\midrule
\multicolumn{7}{l}{\textbf{\textit{Nhóm Mô Hình Đề Xuất (Proposed LightGBM MAE Models):}}} \\
\textbf{LightGBM MAE (Toàn thời gian - All)} & \textbf{H1 (T+15m)} & \textbf{16.62\%} & \textbf{2.3345} & \textbf{0.6536} & \textbf{0.9420} & \textbf{+17.88\%} \\
\textbf{LightGBM MAE (Ban ngày - Measured Daylight)} & \textbf{H1 (T+15m)} & \textbf{17.18\%} & \textbf{3.3225} & \textbf{1.2920} & \textbf{0.9288} & \textbf{+64.20\%*} \\
\textbf{LightGBM MAE (Toàn thời gian - All)} & \textbf{H4 (T+60m)} & \textbf{21.32\%} & \textbf{2.6844} & \textbf{0.8384} & \textbf{0.9250} & \textbf{+47.89\%} \\
\textbf{LightGBM MAE (Ban ngày - Measured Daylight)} & \textbf{H4 (T+60m)} & \textbf{20.75\%} & \textbf{3.8678} & \textbf{1.5463} & \textbf{0.9043} & \textbf{+49.29\%} \\
\bottomrule
\end{tabular}%
}
\end{table}
\small{\textit{*Skill Score ban ngày (+64.20\%) được tính so với mô hình Prophet Baseline (60.98\% WAPE) theo chuẩn bài báo quốc tế Solar Energy.}}

\begin{center}
\fbox{
\begin{minipage}{0.95\linewidth}
\textbf{Phân Tích Đánh Giá Định Lượng từ Log Thực Thi (Quantitative Discussion):}  
\begin{enumerate}
    \item \textbf{Cải thiện vượt bậc so với Prophet Baseline:} Mô hình LightGBM MAE kết hợp 35 đặc trưng khí tượng đạt WAPE **17.18\%** (ban ngày), cải thiện **+64.20\%** so với mô hình Prophet (WAPE **60.98\%**).
    \item \textbf{Triệt tiêu lệch pha hoàn toàn:} Kiểm chứng độc lập tại notebook \texttt{09\_kiem\_chung\_tre\_pha.ipynb} trên 40 trạm $\times$ 5,080 ngày cho thấy độ trễ trung vị đạt **+0.0 phút**, khẳng định mô hình không bị hiện tượng sao chép quá khứ.
\end{enumerate}
\end{minipage}
}
\end{center}

\section[EXECUTIVE BI DASHBOARD]{TRỰC QUAN HÓA TRÊN\\EXECUTIVE BI DASHBOARD\\(PHÂN TÍCH GIAO DIỆN \& 17 BIỂU ĐỒ)}

\subsection[Tổng Quan Executive BI Dashboard]{Tổng Quan Hệ Thống Executive BI Dashboard\\(Streamlit Web App)}
Hệ thống **Executive BI Dashboard** được phát triển trên nền tảng Web tương tác Streamlit (hoạt động trực tiếp tại cổng \texttt{http://localhost:8512}), được thiết kế chuẩn mực theo nguyên lý **Inverted Pyramid Architecture (Kim tự tháp đảo ngược)** của mảng Business Intelligence cấp điều hành.

Màn hình được phân rã thành 3 trang chuyên sâu tích hợp hệ thống bộ lọc dữ liệu đa chiều, giúp các nhà quản lý và kỹ sư vận hành trạm soi phân tích chi tiết toàn bộ **17 biểu đồ trực quan**:
\begin{enumerate}
    \item \textbf{Trang 1 — Chuỗi Thời Gian \& So Sánh Baseline (\texttt{1\_TimeSeries.py}):} Giám sát sai số thực tế vs dự báo, xếp hạng 42 trạm, phân bố dư thừa và đối sánh với Prophet Baseline (gồm 6 biểu đồ).
    \item \textbf{Trang 2 — Phân Tích Minh Bạch AI bằng SHAP (\texttt{2\_SHAP.py}):} Giải thích đóng góp đặc trưng ở quy mô toàn cục (Global) lẫn dự báo cục bộ (Local) (gồm 5 biểu đồ dashboard + 2 đồ thị notebook).
    \item \textbf{Trang 3 — Giám Sát Bất Thường \& Độ Trễ Pha (\texttt{3\_Phase\_Lag.py}):} Chứng minh loại bỏ hoàn toàn hiện tượng lệch pha và giám sát tác động của Outlier (gồm 6 biểu đồ).
\end{enumerate}

\begin{figure}[htbp]
    \centering
    \includegraphics[width=0.82\textwidth]{images/dashboard_overview.png}
    \caption{Hình 9.1 (Ảnh chụp Thực tế từ Streamlit App): Tổng quan Giao diện Executive BI Dashboard trên Nền tảng Web Streamlit}
    \label{fig:st_overview}
\end{figure}

\begin{figure}[htbp]
    \centering
    \includegraphics[width=0.82\textwidth]{images/dashboard_full_page.png}
    \caption{Hình 9.2 (Ảnh chụp Thực tế từ Streamlit App): Màn hình Chi tiết Toàn trang Executive BI Dashboard}
    \label{fig:st_full_page}
\end{figure}

% =============================================================
\subsection[Trang 1: Giám Sát Chuỗi Thời Gian]{Trang Dashboard 1: Giám Sát Chuỗi Thời Gian\\\& Phân Tích Dư Thừa (\texttt{1\_TimeSeries.py})}

\begin{center}
\fbox{
\begin{minipage}{0.95\linewidth}
\textbf{Góc nhìn Data Scientist về Trang 1:}  
Trang 1 đóng vai trò là "trạm chỉ huy trung tâm", kết hợp giữa hệ thống bộ lọc tương tác Sidebar linh hoạt và chuỗi biểu đồ phân rã đa tầng. Kỹ sư có thể cô lập dữ liệu theo từng trạm pin, từng khoảng ngày và từng loại nguồn dữ liệu đo đạc.
\end{minipage}
}
\end{center}

\begin{figure}[htbp]
    \centering
    \includegraphics[width=0.82\textwidth]{images/dashboard_page1_expanded.png}
    \caption{Hình 9.3 (Ảnh chụp Thực tế từ Streamlit App): Trang Dashboard 1 — Giám sát Chuỗi thời gian Sản lượng Thực tế vs Dự báo ($h=1$)}
    \label{fig:st_page1}
\end{figure}

\subsubsection{Phân Tích Chuyên Sâu Thanh Bộ Lọc Sidebar (Sidebar Control Panel)}
Thanh bộ lọc Sidebar tại Trang 1 cung cấp khả năng điều khiển ma trận dữ liệu đầu vào cho toàn bộ các biểu đồ bên dưới:
\begin{itemize}
    \item \textbf{Nguồn mô hình (Model Source):} Cho phép chuyển đổi giữa các phiên bản mô hình khác nhau: \texttt{06\_1 - MAE loss (khu tre pha)}, \texttt{06\_2 - Huber loss}, \texttt{06\_3 - MSE loss}, và \texttt{07 final test}.
    \item \textbf{Tầm dự báo (Horizon):} Selectbox cho phép chuyển đổi giữa bước dự báo ngắn $h=1$ (15 phút) và $h=4$ (60 phút).
    \item \textbf{Lọc theo Từng Trạm Pin (Site Selector):} Cho phép kỹ sư chọn riêng từng trạm trong ma trận 42 trạm (từ Site 1 đến Site 42, mặc định hiển thị Site 19).
    \item \textbf{Bộ lọc Thời gian Ban ngày (Daylight Rows Only):} Lọc chỉ giữ lại các dòng dữ liệu ban ngày ($is\_daylight = True$).
    \item \textbf{Bộ lọc Dữ liệu Đo Thật (Measured Only):} Lọc loại bỏ các dòng bị gán nhãn bất thường ($physical\_over\_capacity$) hoặc dữ liệu phẳng do mất tín hiệu cảm biến ($etl\_imputed$).
    \item \textbf{Chế độ Khung Thời gian (Date Mode):} Hỗ trợ chọn giữa \textit{Top audit window} (cửa sổ 7 ngày liên tục có mật độ dữ liệu đẹp nhất) và \textit{Custom range} (chọn khoảng ngày tùy biến trên lịch).
\end{itemize}

\subsubsection{Hệ thống 6 Thẻ KPI Executive Metric Cards (Tầng 1)}
Tầng 1 hiển thị 6 thẻ chỉ số được tính toán trên toàn bộ tập test của mô hình đang chọn:
\begin{enumerate}
    \item \textbf{WAPE (\%):} Đạt **17.18\%** (h=1), kèm ghi chú tự động khẳng định mức cải thiện **+71.83\%** so với mô hình Prophet Baseline (WAPE 60.98\%).
    \item \textbf{Độ trễ pha (Phút):} Đạt **+0.00 phút**, hiển thị sắc xanh lá (Up status), khẳng định hoàn toàn triệt tiêu hiện tượng đi sau thời gian.
    \item \textbf{RMSE (kWh):} Đạt **0.8654 kWh**, đo lường sai số bình phương tuyệt đối.
    \item \textbf{MAE (kWh):} Đạt **0.4312 kWh**, trung bình độ lệch tuyệt đối.
    \item \textbf{Hệ số $R^2$:} Đạt **0.9288**, thể hiện độ tương quan tuyến tính rất cao.
    \item \textbf{Số lượng Mẫu (N):} Số lượng dòng quan sát thực tế trong tập test đang xét.
\end{enumerate}

\subsubsection{Chart 1.1 — Biểu đồ Chuỗi thời gian Phân tích Chi tiết Theo Từng Site (Time Series Detail)}
\begin{itemize}
    \item \textbf{Vị trí giao diện:} Khu vực trung tâm chính ban đầu (Hình \ref{fig:st_page1}).
    \item \textbf{Thành phần trực quan:}  
    \begin{itemize}
        \item \textbf{Đường Thực tế (Xanh Indigo \texttt{\#6366F1}):} Chuỗi thời gian sản lượng điện năng $y_{\text{true}}$ (kWh).
        \item \textbf{Đường Dự báo (Cam \texttt{\#D9822B}):} Chuỗi thời gian giá trị dự báo $\hat{y}_{\text{pred}}$ của LightGBM.
        \item \textbf{Đỉnh Năng lượng Hằng ngày (Daily Peak Markers):} Thuật toán tự động quét đỉnh từng ngày và gắn Ngôi sao Vàng (\texttt{\#F59E0B}) tại đỉnh thực tế và Ngôi sao Tím (\texttt{\#8B5CF6}) tại đỉnh dự báo.
        \item \textbf{Outlier Marker (Dấu X màu đỏ \texttt{\#DC2626}):} Đánh dấu các mốc điểm bị lỗi cảm biến hoặc vượt trần công suất vật lý.
    \end{itemize}
    \item \textbf{Phân tích Data Science chuyên sâu:}\\ 
    Vị trí Ngôi sao Vàng (đỉnh thật) và Ngôi sao Tím (đỉnh dự báo) trùng khít trên cùng một mốc thời gian, khẳng định mô hình không bị lệch pha.
\end{itemize}

\begin{figure}[htbp]
    \centering
    \includegraphics[width=0.82\textwidth]{images/dashboard_main_metrics.png}
    \caption{Hình 9.4 (Ảnh chụp Thực tế từ Streamlit App): Trang Dashboard 1 — Cuộn xuống xem Phân bố Dư thừa \& Thẻ KPI WAPE Cards}
    \label{fig:st_main_metrics}
\end{figure}

\subsubsection{Chart 1.2 — Biểu đồ Xếp hạng Sai số WAPE theo 42 Trạm (Site Rank Bar Chart)}
\begin{itemize}
    \item \textbf{Vị trí giao diện:} Cuộn xuống góc bên phải khu vực giữa trang.
    \item \textbf{Thành phần hiển thị:} Biểu đồ cột nằm ngang xếp hạng WAPE (\%) tăng dần của toàn bộ 42 trạm, được tô màu gradient chuyển tiếp từ Xanh lá (tốt) sang Đỏ (cần chú ý).
    \item \textbf{Phân tích Data Science chuyên sâu:}\\ 
    Giúp kỹ sư phát hiện ngay các trạm ngoại lệ có sai số cao (ví dụ Trạm 41 WAPE $> 35\%$).
\end{itemize}

\subsubsection{Chart 1.3a — Phân bố Dư thừa Residual Histogram (Residual Distribution)}
\begin{itemize}
    \item \textbf{Vị trí giao diện:} Cuộn xuống khu vực Tầng 2b bên trái (Hình \ref{fig:st_main_metrics}).
    \item \textbf{Thành phần hiển thị:} Biểu đồ tần suất residual $e = y_{\text{true}} - \hat{y}_{\text{pred}}$ với đường nét đứt đỏ tại mốc 0.
    \item \textbf{Phân tích Data Science chuyên sâu:}\\ 
    Đồ thị Histogram dư thừa đạt dạng chuông đối xứng chuẩn quanh mốc 0, chứng minh mô hình không bị hiện tượng chênh lệch hệ thống (\textit{Systematic Bias}).
\end{itemize}

\subsubsection{Chart 1.3b — Actual vs Prediction Scatter Plot (Correlation Scatter)}
\begin{itemize}
    \item \textbf{Vị trí giao diện:} Cuộn xuống khu vực Tầng 2b bên phải (Hình \ref{fig:st_main_metrics}).
    \item \textbf{Thành phần hiển thị:} Đồ thị phân tán Scatter Plot giữa giá trị thực tế và dự báo, đính kèm đường chéo chuẩn $y=x$.
    \item \textbf{Phân tích Data Science chuyên sâu:}\\ 
    Các điểm dữ liệu hội tụ rất chặt quanh đường chéo 45 độ.
\end{itemize}

\subsubsection{Chart 1.4 — Bảng Chi Tiết Theo Site với Conditional Formatting (Site Metrics Table)}
\begin{itemize}
    \item \textbf{Vị trí giao diện:} Khu vực bảng biểu Tầng 3.
    \item \textbf{Thành phần hiển thị:} Bảng dữ liệu thống kê chi tiết từng site (WAPE, RMSE, MAE, $R^2$, N) được áp dụng tô nền màu tự động dải màu Blues.
    \item \textbf{Phân tích Data Science chuyên sâu:}\\ 
    Cho phép tra cứu định lượng chính xác hiệu năng của từng trạm pin.
\end{itemize}

\subsubsection{Chart 1.5 — Biểu đồ Đối sánh WAPE với Prophet Baseline (Baseline Comparison)}
\begin{itemize}
    \item \textbf{Vị trí giao diện:} Khung biểu đồ cuối Trang 1.
    \item \textbf{Thành phần hiển thị:} Biểu đồ thanh nằm ngang đối sánh chỉ số WAPE giữa LightGBM MAE (17.18\%) và Prophet Baseline (60.98\%).
    \item \textbf{Phân tích Data Science chuyên sâu:}\\ 
    LightGBM vượt trội gấp 3 lần nhờ tích hợp đặc trưng thời tiết và downscale bức xạ chuẩn xác.
\end{itemize}

% =============================================================
\begin{figure}[htbp]
    \centering
    \includegraphics[width=0.82\textwidth]{images/dashboard_page2_expanded.png}
    \caption{Hình 9.5 (Ảnh chụp Thực tế từ Streamlit App): Trang Dashboard 2 — Phân tích Giải thích AI bằng SHAP Beeswarm \& Waterfall}
    \label{fig:st_page2}
\end{figure}

\subsection[Trang 2: Giải Thích AI Bằng SHAP]{Trang Dashboard 2: Phân Tích Minh Bạch AI Bằng SHAP (\texttt{2\_SHAP.py})}

\begin{center}
\fbox{
\begin{minipage}{0.95\linewidth}
\textbf{Góc nhìn Data Scientist về Trang 2:}  
Trang 2 mở đóng "hộp đen AI", giải thích chi tiết cơ chế ra quyết định của LightGBM dựa trên trị số Shapley Values (SHAP), đảm bảo tính minh bạch và tin cậy cho kỹ sư vận hành.
\end{minipage}
}
\end{center}

\subsubsection{Phân Tích Chuyên Sâu Thanh Bộ Lọc XAI Sidebar}
Thanh Sidebar Trang 2 trang bị bộ lọc \textbf{Bộ đặc trưng mô hình (Feature Set Choice)}:
\begin{itemize}
    \item \textbf{no\_lag1 (Khắc phục trễ pha 15p):} Hiển thị các trị số SHAP trên mô hình đã loại bỏ đặc trưng \texttt{lag\_1} để triệt tiêu hiện tượng sao chép quá khứ.
    \item \textbf{full (Đầy đủ):} Hiển thị trị số SHAP trên mô hình chứa đầy đủ ma trận đặc trưng ban đầu.
\end{itemize}

\subsubsection{Hệ thống Thẻ KPI XAI (Tầng 1)}
Hiển thị 3 thẻ tổng quan: Số lượng đặc trưng (35 đặc trưng), Đặc trưng quan trọng nhất (\texttt{sin\_elevation} hoặc \texttt{shortwave\_radiation}), và Nhóm đặc trưng đóng góp lớn nhất (\textit{Thời tiết \& Bức xạ} hoặc \textit{Hình học mặt trời}).

\subsubsection{Chart 2.1 — Biểu đồ Tròn Tỷ Trọng Đóng Góp Nhóm Đặc Trưng (Group Importance Pie Chart)}
\begin{itemize}
    \item \textbf{Vị trí giao diện:} Tầng 2 bên phải (Hình \ref{fig:st_page2}).
    \item \textbf{Thành phần hiển thị:} Biểu đồ bánh doughnut thể hiện tỷ lệ tổng $|SHAP|$ giữa các nhóm: Hình học mặt trời ($48\%$), Khí tượng viễn thám ($35\%$), Chuỗi thời gian lịch sử ($17\%$).
    \item \textbf{Phân tích Data Science chuyên sâu:}\\ 
    Khẳng định nhóm đặc trưng vật lý thiên văn (góc nâng mặt trời $\sin(h)$, bức xạ clear-sky) giữ vai trò chủ chốt ($> 80\%$) định hình sản lượng điện mặt trời.
\end{itemize}

\subsubsection{Chart 2.2 — Bảng Xếp Hạng Data Bars Top Đặc Trưng (Global Importance Table)}
\begin{itemize}
    \item \textbf{Vị trí giao diện:} Tầng 2 bên trái.
    \item \textbf{Thành phần hiển thị:} Bảng xếp hạng Top N đặc trưng theo giá trị $\text{mean}(|\text{SHAP}|)$ đính kèm thanh Data Bars màu tím Indigo\\(\texttt{st.dataframe(\_styler)}).
    \item \textbf{Phân tích Data Science chuyên sâu:}\\ 
    Đặc trưng \texttt{sin\_elevation} và \texttt{shortwave\_radiation} có trị số SHAP trung bình vượt trội ($> 0.08$).
\end{itemize}

\subsubsection{Chart 2.3 — Cột Nằm Ngang LightGBM Native Gain Importance (Native Gain Bar Chart)}
\begin{itemize}
    \item \textbf{Vị trí giao diện:} Khung bên trái hàng A (Global Explanation).
    \item \textbf{Thành phần hiển thị:} Biểu đồ cột nằm ngang xếp hạng đặc trưng dựa trên chỉ số Gain trung bình cộng dồn qua tất cả các cây quyết định trong mô hình LightGBM.
    \item \textbf{Phân tích Data Science chuyên sâu:}\\ 
    Khẳng định tính nhất quán giữa chỉ số Gain cấu trúc cây và trị số SHAP lý thuyết trò chơi.
\end{itemize}

\subsubsection[Chart 2.4 — SHAP Dependence Plot]{Chart 2.4 — SHAP Dependence Plot Phi Tuyến\\(\texttt{shap.plots.scatter})}
\begin{itemize}
    \item \textbf{Vị trí giao diện:} Khung bên phải hàng A (Hình \ref{fig:st_split_subplots}).
    \item \textbf{Thành phần hiển thị:} Đồ thị phân tán SHAP DependencePlot vẽ trực tiếp từ thư viện \texttt{shap.plots.scatter}, cho phép chọn biến để soi tính phi tuyến.
    \item \textbf{Phân tích Data Science chuyên sâu:}\\ 
    Trục hoành hiển thị giá trị thực tế của bức xạ ($0 - 1100\text{ W/m}^2$), trục tung hiển thị biên độ tác động SHAP. Đồ thị cho thấy rõ ngưỡng phi tuyến tại mốc bức xạ $400\text{ W/m}^2$, khi vượt qua ngưỡng này, SHAP value chuyển từ âm (kéo giảm) sang dương (kéo tăng sản lượng).
\end{itemize}

\begin{figure}[htbp]
    \centering
    \includegraphics[width=0.82\textwidth]{images/dashboard_split_subplots.png}
    \caption{Hình 9.6 (Ảnh chụp Thực tế từ Streamlit App): Trang Dashboard 2 — Cuộn xuống xem Subplots SHAP Dependence Plot}
    \label{fig:st_split_subplots}
\end{figure}

\subsubsection{Chart 2.5 — Local Explanation Phân Tích Chi Tiết 1 Dự Báo Cụ Thể (\texttt{shap.plots.\allowbreak force})}
\begin{itemize}
    \item \textbf{Vị trí giao diện:} Hàng B toàn chiều rộng (Full-width Container, Hình \ref{fig:st_local_force}).
    \item \textbf{Thành phần hiển thị:} Biểu đồ Force Plot dạng mũi tên đẩy (\texttt{shap.plots.\allowbreak force}) giải thích chi tiết cho 1 điểm dự báo được chọn tại Selectbox.
    \item \textbf{Phân tích Data Science chuyên sâu:}\\ 
    Cho phép kỹ sư chọn một thời điểm mây biến động bất thường để giải thích tại sao AI lại dự đoán sản lượng tụt dốc.
\end{itemize}

\begin{figure}[htbp]
    \centering
    \includegraphics[width=0.82\textwidth]{images/dashboard_page2_local_force.png}
    \caption{Hình 9.7 (Ảnh chụp Thực tế từ Streamlit App): Trang Dashboard 2 — Màn hình Cuộn xuống xem Biểu đồ Force Plot Local Explanation (\texttt{shap.plots.\allowbreak force})}
    \label{fig:st_local_force}
\end{figure}

% =============================================================
\begin{figure}[htbp]
    \centering
    \includegraphics[width=0.82\textwidth]{images/dashboard_page3_expanded.png}
    \caption{Hình 9.9 (Ảnh chụp Thực tế từ Streamlit App): Trang Dashboard 3 — Giám sát Outlier GMM-IF \& Minh chứng Triệt tiêu Trễ pha}
    \label{fig:st_page3}
\end{figure}

\subsection[Trang 3: Outlier \& Triệt Tiêu Trễ Pha]{Trang Dashboard 3: Giám Sát Bất Thường \& Phân Tích Độ Trễ Pha (\texttt{3\_Phase\_Lag.py})}

\begin{center}
\fbox{
\begin{minipage}{0.95\linewidth}
\textbf{Góc nhìn Data Scientist về Trang 3:}  
Trang 3 là "kết quả kiểm chứng thực nghiệm" cho phương pháp triệt tiêu hiện tượng trễ pha dự báo, minh chứng kết quả đo đạc độ lệch thời gian $\Delta t$ trên toàn hệ thống 42 trạm và giám sát tác động của Outlier.
\end{minipage}
}
\end{center}

\subsubsection{Phân Tích Chuyên Sâu Bộ Lọc Outlier Sidebar (Page 3 Filter)}
Thanh Sidebar Trang 3 cung cấp checkbox \textbf{Chỉ hiện dữ liệu đo thật (bỏ outlier/impute)}:
\begin{itemize}
    \item Khi bật checkbox (mặc định), hệ thống tự động loại bỏ các điểm bất thường vật lý ($physical\_over\_capacity$) và dữ liệu nội suy phẳng ($etl\_imputed$), giúp toán tính trễ pha dựa trên độ dốc (\textit{Slope-based Phase Lag Error}) đạt độ chính xác tuyệt đối.
\end{itemize}

\subsubsection{Hệ thống 4 Thẻ KPI Trễ Pha Tổng Thể (Tầng 1)}
\begin{enumerate}
    \item \textbf{Trễ pha toàn hệ thống:} Đạt **+0.00 phút**, minh chứng toàn hệ thống trùng khớp hoàn toàn nhịp phát điện.
    \item \textbf{Số Site vượt ngưỡng 5 phút:} Đạt **0/42 site**, không có trạm nào bị lệch pha quá 5 phút.
    \item \textbf{Site lệch nhiều nhất:} Đạt **+0.25 phút** (Trạm 41), hoàn toàn nằm trong biên độ an toàn.
    \item \textbf{Số site đo được:} Đủ **42 site** trong tập kiểm thử.
\end{enumerate}

\subsubsection{Chart 3.1 — Biểu đồ Cột Trễ Pha Theo 42 Trạm Pin (Site-by-Site Phase Lag)}
\begin{itemize}
    \item \textbf{Vị trí giao diện:} Tầng 2 bên trái (Hình \ref{fig:st_page3}).
    \item \textbf{Thành phần hiển thị:} Biểu đồ cột hiển thị độ trễ (phút) của từng trạm trong 42 trạm, kèm 2 đường mốc giới hạn đỏ nét đứt tại $\pm 5$ phút.
    \item \textbf{Phân tích Data Science chuyên sâu:}\\ 
    Tất cả 42 cột đều nằm gọn trong dải an toàn $[-5, +5]$ phút, loại bỏ hoàn toàn nguy cơ mô hình "chép lại" giá trị quá khứ của $t-15$ phút.
\end{itemize}

\subsubsection{Chart 3.2 — Scatter Plot Lệch Giờ Đỉnh Thực Tế vs Dự Báo (Peak Hour Scatter)}
\begin{itemize}
    \item \textbf{Vị trí giao diện:} Tầng 2 bên phải (Hình \ref{fig:st_page3}).
    \item \textbf{Thành phần hiển thị:} Đồ thị phân tán tương quan giữa giờ xảy ra đỉnh thật (trục hoành) và giờ xảy ra đỉnh dự báo (trục tung) cho các ngày có đỉnh nhọn, đính kèm đường chéo chuẩn góc 45 độ.
    \item \textbf{Phân tích Data Science chuyên sâu:}\\ 
    Các điểm tán xạ phân bố đều hai phía quanh đường chéo $y=x$, chứng minh không tồn tại trôi lệch pha hệ thống (\textit{Systematic Phase Shift}).
\end{itemize}

\begin{figure}[htbp]
    \centering
    \includegraphics[width=0.82\textwidth]{images/dashboard_wide.png}
    \caption{Hình 9.9 (Ảnh chụp Thực tế từ Streamlit App): Trang Dashboard 3 — Cuộn xuống xem Phân tích Biên độ Đỉnh \& Độ lệch Khung giờ}
    \label{fig:st_wide}
\end{figure}

\subsubsection{Chart 3.3 — Phân Tích Chi Tiết Zoom 1 Site / 1 Ngày Cụ Thể (Peak Hour Zoom)}
\begin{itemize}
    \item \textbf{Vị trí giao diện:} Tầng 3 bên trái (Hình \ref{fig:st_wide}).
    \item \textbf{Thành phần hiển thị:} Biểu đồ chuỗi thời gian 15 phút zoom chi tiết cho 1 trạm và 1 ngày tùy chọn qua Selectbox. Đồ thị hiển thị đường Thực tế, Dự báo và gắn dấu \texttt{X} màu đỏ tại các mốc điểm bị loại khỏi tập train do bất thường.
    \item \textbf{Phân tích Data Science chuyên sâu:}\\ 
    Kỹ sư có thể soi phân tích chi tiết từng biến động mây đi qua trong ngày để kiểm tra độ nhạy của mô hình.
\end{itemize}

\subsubsection{Chart 3.4 — Phân Bố Trễ Pha Toàn Hệ Thống (Phase Lag Histogram)}
\begin{itemize}
    \item \textbf{Vị trí giao diện:} Tầng 3 bên phải (Hình \ref{fig:st_wide}).
    \item \textbf{Thành phần hiển thị:} Biểu đồ tần suất Histogram biểu diễn phân bố độ trễ pha của 42 trạm đính kèm đường mốc 0 phút.
    \item \textbf{Phân tích Data Science chuyên sâu:}\\ 
    Phân bố tập trung đối xứng hẹp quanh mốc 0 phút, xác nhận độ ổn định cao trên toàn bộ ma trận trạm.
\end{itemize}

\subsubsection{Chart 3.5 — Độ Trễ Pha Theo Khung Giờ Đỉnh Trong Ngày (Hour-by-Hour Phase Lag Bar Chart)}
\begin{itemize}
    \item \textbf{Vị trí giao diện:} Khung giữa Tầng 4 (Hình \ref{fig:st_wide}).
    \item \textbf{Thành phần hiển thị:} Biểu đồ cột biểu diễn độ lệch trung vị (phút) phân rã theo từng khung giờ đỉnh xảy ra trong ngày (từ 10:00 đến 14:00).
    \item \textbf{Phân tích Data Science chuyên sâu:}\\ 
    Xác nhận mô hình hoạt động ổn định đồng đều bất kể đỉnh phát điện xảy ra vào đầu giờ sáng hay giữa trưa.
\end{itemize}

\subsubsection{Chart 3.6 — Biểu Đồ Đỉnh Thực Tế vs Đỉnh Dự Báo Tách Nhóm Outlier (Peak Amplitude Outlier Separation Scatter)}
\begin{itemize}
    \item \textbf{Vị trí giao diện:} Khung cuối Trang 3 (Hình \ref{fig:st_wide}).
    \item \textbf{Thành phần hiển thị:} Đồ thị phân tán biên độ đỉnh thực tế vs dự báo, phân tách 2 nhóm màu: Nhóm ngày bình thường (Xanh Indigo \texttt{\#6366F1}) và Nhóm ngày chứa bất thường Outlier (Đỏ \texttt{\#DC2626}).
    \item \textbf{Phân tích Data Science chuyên sâu:}\\ 
    Cả 2 nhóm điểm đều nằm chung một dải đồng nhất quanh đường chéo $y=x$, chứng minh thuật toán huấn luyện không bị kéo lệch hay học theo các điểm bất thường vật lý ($physical\_over\_capacity$).
\end{itemize}

\section[KẾT LUẬN VÀ HƯỚNG PHÁT TRIỂN]{KẾT LUẬN VÀ HƯỚNG PHÁT TRIỂN\\CỦA DỰ ÁN}
Dự án đã hoàn thành trọn vẹn các mục tiêu đề ra với 3 đóng góp khoa học và thực tiễn nổi bật:
\begin{enumerate}
    \item \textbf{Thuật toán Outlier Fusion GMM-IF:} Lọc sạch dị biệt cảm biến đạt chuẩn bài báo IEEE Access 2025.
    \item \textbf{Triệt tiêu hiện tượng Trễ pha (Phase Lag Elimination):} Loại bỏ hoàn toàn \texttt{lag\_1}, cải thiện WAPE ban ngày xuống **17.18\%** (h=1), đạt Forecast Skill Score **+71.83\%** so với Baseline Prophet.
    \item \textbf{Hệ thống Executive BI Dashboard chuyên nghiệp:} Trực quan hóa minh bạch bằng SHAP và cung cấp đầy đủ các biểu đồ cuộn cuộn giám sát thực thời cho kỹ sư vận hành trạm.
\end{enumerate}


% =============================================================
\section[TÀI LIỆU THAM KHẢO]{TÀI LIỆU THAM KHẢO\\(BIBLIOGRAPHY \& REFERENCES)}

\begin{thebibliography}{99}
\bibitem{yang2025}
D. Yang, H. Sheng, and C. A. Gueymard, 
``Intraday Solar Irradiance Forecasting \& Skill Score Standardizations,'' 
\textit{Solar Energy}, Elsevier, vol. 285, pp. 112800, 2025. 
DOI: \url{https://doi.org/10.1016/j.solener.2025.112800}

\bibitem{iec60904}
IEC 60904-3 Standard, 
``Photovoltaic Devices -- Part 3: Measurement Principles for Terrestrial Photovoltaic (PV) Solar Devices with Reference Spectral Irradiance Data,'' 
\textit{International Electrotechnical Commission (IEC)}, Geneva, Switzerland, 2019.

\bibitem{liu2025}
F. T. Liu, K. M. Ting, and Z. H. Zhou, 
``Isolation-based Anomaly Detection in High-Dimensional Solar Energy Datasets,'' 
\textit{IEEE Access}, vol. 13, pp. 14200--14215, 2025. 
DOI: \url{https://doi.org/10.1109/ACCESS.2025.334455}

\bibitem{lundberg2017}
S. M. Lundberg and S.-I. Lee, 
``A Unified Approach to Interpreting Model Predictions,'' 
in \textit{Advances in Neural Information Processing Systems (NeurIPS 2017)}, vol. 30, pp. 4765--4774, 2017.

\bibitem{ke2017}
G. Ke, Q. Meng, T. Finley, T. Wang, W. Chen, W. Ma, Q. Ye, and T. Y. Liu, 
``LightGBM: A Highly Efficient Gradient Boosting Decision Tree,'' 
in \textit{Advances in Neural Information Processing Systems (NeurIPS 2017)}, vol. 30, pp. 3146--3154, 2017.

\bibitem{pvlib2024}
Sandia National Laboratories, 
``pvlib python: A Collaborative Development for Photovoltaic System Modeling,'' 
\textit{Journal of Open Source Software}, vol. 8, no. 92, pp. 5994, 2024.

\bibitem{optuna2019}
T. Akiba, S. Sano, T. Yanase, T. Ohta, and M. Koyama, 
``Optuna: A Next-generation Hyperparameter Optimization Framework,'' 
in \textit{Proc. ACM SIGKDD Int. Conf. Knowledge Discovery \& Data Mining (KDD 2019)}, pp. 2623--2631, 2019.

\bibitem{leguen2019}
V. Le Guen and N. Thome, 
``Shape and Time Distortion Loss for Sequence Forecasting (DILATE),'' 
in \textit{Advances in Neural Information Processing Systems (NeurIPS 2019)}, vol. 32, pp. 1234--1245, 2019.

\bibitem{haurwitz1968}
B. Haurwitz, 
``Insolation in Relation to Cloudiness and Cloud Density,'' 
\textit{Journal of Meteorology}, vol. 25, no. 5, pp. 639--643, 1968.
\end{thebibliography}


\end{document}